% Options for packages loaded elsewhere
\PassOptionsToPackage{unicode}{hyperref}
\PassOptionsToPackage{hyphens}{url}
\documentclass[
]{article}
\usepackage{xcolor}
\usepackage{amsmath,amssymb}
\setcounter{secnumdepth}{-\maxdimen} % remove section numbering
\usepackage{iftex}
\ifPDFTeX
  \usepackage[T1]{fontenc}
  \usepackage[utf8]{inputenc}
  \usepackage{textcomp} % provide euro and other symbols
\else % if luatex or xetex
  \usepackage{unicode-math} % this also loads fontspec
  \defaultfontfeatures{Scale=MatchLowercase}
  \defaultfontfeatures[\rmfamily]{Ligatures=TeX,Scale=1}
\fi
\usepackage{lmodern}
\ifPDFTeX\else
  % xetex/luatex font selection
\fi
% Use upquote if available, for straight quotes in verbatim environments
\IfFileExists{upquote.sty}{\usepackage{upquote}}{}
\IfFileExists{microtype.sty}{% use microtype if available
  \usepackage[]{microtype}
  \UseMicrotypeSet[protrusion]{basicmath} % disable protrusion for tt fonts
}{}
\makeatletter
\@ifundefined{KOMAClassName}{% if non-KOMA class
  \IfFileExists{parskip.sty}{%
    \usepackage{parskip}
  }{% else
    \setlength{\parindent}{0pt}
    \setlength{\parskip}{6pt plus 2pt minus 1pt}}
}{% if KOMA class
  \KOMAoptions{parskip=half}}
\makeatother
\usepackage{longtable,booktabs,array}
\usepackage{calc} % for calculating minipage widths
% Correct order of tables after \paragraph or \subparagraph
\usepackage{etoolbox}
\makeatletter
\patchcmd\longtable{\par}{\if@noskipsec\mbox{}\fi\par}{}{}
\makeatother
% Allow footnotes in longtable head/foot
\IfFileExists{footnotehyper.sty}{\usepackage{footnotehyper}}{\usepackage{footnote}}
\makesavenoteenv{longtable}
\setlength{\emergencystretch}{3em} % prevent overfull lines
\providecommand{\tightlist}{%
  \setlength{\itemsep}{0pt}\setlength{\parskip}{0pt}}
\usepackage{microtype}
\usepackage{titlesec}
\usepackage{booktabs}
\usepackage{array}
\usepackage{longtable}
\usepackage{ragged2e}
\usepackage{etoolbox}
\usepackage{caption}
\usepackage{fancyhdr}
\usepackage{setspace}
\usepackage{newunicodechar}

\newunicodechar{✓}{$\checkmark$}
\newunicodechar{∈}{$\in$}
\newunicodechar{∃}{$\exists$}
\newunicodechar{≠}{$\neq$}
\newunicodechar{≤}{$\le$}
\newunicodechar{≥}{$\ge$}
\newunicodechar{→}{$\rightarrow$}
\newunicodechar{₁}{$_1$}
\newunicodechar{₂}{$_2$}
\newunicodechar{ₖ}{$_k$}
\newunicodechar{─}{-}
\newunicodechar{└}{+}
\newunicodechar{├}{+}
\newunicodechar{│}{|}

\setstretch{1.03}
\setlength{\parindent}{0pt}
\setlength{\parskip}{0.45em}
\setlength{\emergencystretch}{3em}

\titleformat{\section}{\Large\bfseries}{\thesection.}{0.6em}{}
\titleformat{\subsection}{\large\bfseries}{\thesubsection.}{0.6em}{}
\titleformat{\subsubsection}{\normalsize\bfseries}{\thesubsubsection.}{0.6em}{}

\titlespacing*{\section}{0pt}{1.4em}{0.6em}
\titlespacing*{\subsection}{0pt}{1.1em}{0.4em}
\titlespacing*{\subsubsection}{0pt}{0.8em}{0.3em}

\setlength{\LTpre}{0.35em}
\setlength{\LTpost}{0.35em}
\setlength{\LTleft}{0pt}
\setlength{\LTright}{0pt}
\setlength{\tabcolsep}{4pt}
\renewcommand{\arraystretch}{1.08}
\AtBeginEnvironment{longtable}{\small}
\newcolumntype{L}[1]{>{\RaggedRight\arraybackslash}p{#1}}
\newcolumntype{C}[1]{>{\Centering\arraybackslash}p{#1}}
\newcolumntype{R}[1]{>{\RaggedLeft\arraybackslash}p{#1}}

\captionsetup[table]{skip=4pt}
\captionsetup[figure]{skip=4pt}

\pagestyle{fancy}
\fancyhf{}
\fancyfoot[C]{\thepage}
\renewcommand{\headrulewidth}{0pt}
\renewcommand{\footrulewidth}{0pt}
\setlength{\headheight}{14pt}

\newcommand{\codeid}[1]{{\small\begingroup\urlstyle{tt}\nolinkurl{#1}\endgroup}}
\usepackage{bookmark}
\IfFileExists{xurl.sty}{\usepackage{xurl}}{} % add URL line breaks if available
\urlstyle{same}
\hypersetup{
  hidelinks,
  pdfcreator={LaTeX via pandoc}}

\author{}
\date{}

\begin{document}

\section{AppSec Core v1: A Formalized Normalization Ontology for
Application
Security}\label{appsec-core-v1-a-formalized-normalization-ontology-for-application-security}

\textbf{Pedro Farinha} Independent Researcher
\href{mailto:pedro.farinha@shiftleft.pt}{\nolinkurl{pedro.farinha@shiftleft.pt}}

\emph{Supported by Shiftleft --- Secure Software Engineering, lda.}

\begin{center}\rule{0.5\linewidth}{0.5pt}\end{center}

\subsection{Abstract}\label{abstract}

This paper publishes \textbf{AppSec Core v1} as a formalized
\textbf{bounded ontology artefact} for application security.
\emph{Bounded} here means: an extraction of the engineering substance of
heterogeneous AppSec and regulatory sources (the
development-and-maintenance practices a software engineer must satisfy
on a pull request, sprint, or release at engagement scale), deliberately
scoped to exclude governance substance (board reporting, vendor
contractual risk management, regulatory examination response, workforce
training programmes) which is routed to a complementary practitioner
Manual surface. v1 is delivered as an OWL 2 DL ontology populated with
259 typed instances across the same ten domain slices introduced by
AppSec Core v0 {[}1{]} (75 ControlObjectives, 69 Practices, 58
Mechanisms, 57 Artifacts), accompanied by a SHACL Core constraint set
published as a complementary apparatus and validated under both the
SHACL Core reference-implementation validator \texttt{pyshacl} and an
in-house bounded-subset validator at parity for the apparatus's six
ontology shapes. SHACL execution against the populated graph yields
\texttt{sh:conforms\ =\ true} with zero \texttt{sh:Violation}. All
artefacts are SHA-256-pinned at the cycle-close release tag chain.

The paper makes three contributions, each demonstrable by inspection of
the published artefact and its validation outputs. \textbf{(1) AppSec
Core v1 as a formalized artefact} --- a citable, machine-validatable OWL
2 DL ontology and SHACL Core constraint set, replacing the YAML-only v0
surface. \textbf{(2) Schema preservation under multi-source pressure}
--- across the v0→v1 transition, the ten-slice partition, the four
populated entity types, the EvidencePattern declarative class, and the
cross-slice relation model are preserved; under expansion of the
underlying source corpus from the five first-wave sources {[}2{]} to
thirty-one sources at cycle close (a 6.2-fold source-count expansion
absorbed by the design science cycle reported in {[}4{]}), the v0→v1
entity-count delta consists exclusively of additive instance population
governed by an explicit protocol. \textbf{(3) The AppSec Core Change
Request (ACR) protocol with its explicit four-condition promotion
threshold} --- the protocol is published with its threshold
(multi-source convergence at ≥5 independent sources from ≥3
organisational authorities; multi-method convergence;
practitioner-Manual content backing; slice fit) and is demonstrated
through symmetric application across four worked decisions: ACR-001
\emph{Secure Configuration Baseline Integrity} (promoted), ACR-002
\emph{Security Requirements Lifecycle Management} (promoted), ACR-003
\emph{Multi-anchor adjacency candidate batch} (not admitted), and
ACR-004 \emph{Output Rendering Safety / Context-Aware Encoding}
(promoted). Three promotions and one non-admission decided under the
same threshold provide initial evidence that the protocol produces
falsifiable governance decisions rather than retrofitting outcomes; the
discriminative power of the threshold against single-condition marginal
failures is established by the threshold's specification, with
empirically grounded near-miss cases registered as future work (§10.6).

The design science cycle that produced v1 from v0, the per-source corpus
analysis, and the bidirectional validation of the practitioner Manual
against v1 are out of scope and are reported in sibling papers {[}4,
5{]}.

\textbf{Keywords:} ontology, application security, OWL 2, SHACL,
normalization, governance protocol, AppSec Core

\begin{center}\rule{0.5\linewidth}{0.5pt}\end{center}

\subsection{1. Introduction}\label{introduction}

\subsubsection{1.1 Background and prior
work}\label{background-and-prior-work}

AppSec Core v0 was introduced in {[}1{]} as a normalized ontology for
application security: ten domain slices with typed instances across four
populated first-class entity types --- ControlObjective, Practice,
Mechanism, Artifact --- plus an EvidencePattern declarative class
indexed for verification semantics. v0 was derived bottom-up from a
practitioner application-security Manual corpus and validated through
canonical reduction of requirements drawn from a deliberately bounded
set of five first-wave external sources: NIST Secure Software
Development Framework (SSDF) {[}9{]}, OWASP Application Security
Verification Standard (ASVS) {[}10{]}, Supply-chain Levels for Software
Artifacts (SLSA) {[}11{]}, CIS Critical Security Controls {[}12{]}, and
MITRE Common Attack Pattern Enumeration and Classification (CAPEC)
{[}13{]}. The v0 paper {[}1{]} established the slice architecture, the
four-type populated instance schema, and the per-slice contracts
specifying scope, non-goals, and acceptance conditions; the companion
compilation paper {[}2{]} applied v0 to a fifteen-chapter practitioner
Manual case study under a coverage-preserving compilation method using a
keyword-first normalization heuristic at cluster-level granularity.

Two limitations of v0 were declared in {[}1{]} and remained open:

\begin{enumerate}
\def\labelenumi{\arabic{enumi}.}
\tightlist
\item
  \textbf{Formalization.} v0 was published as YAML-based typed instance
  schemas with per-slice contracts, sufficient for inspection by human
  readers and for mechanical processing under custom tooling, but not
  directly consumable by standard reasoners or by SHACL validators. The
  v0 paper noted OWL 2 / SHACL formalization as future work.
\item
  \textbf{Source-set bounding.} The five-source first wave was bounded
  by design --- sufficient to demonstrate the ontology's shape and the
  canonical-reduction property, but explicitly insufficient to claim
  that the slice partition would absorb a substantively wider corpus.
  Whether v0 would hold under expanded source pressure was an open
  empirical question.
\end{enumerate}

This paper resolves both limitations. The ontology has been refined
through a design science cycle (reported in {[}4{]}) under expansion of
the underlying source corpus to thirty-one sources drawn from eleven
organisational authorities and including AI-domain references, and the
resulting \textbf{AppSec Core v1} is published here as a formalized OWL
2 DL ontology with a SHACL Core constraint set. The cycle that produced
v1 is not re-derived in this paper; v1 is presented as the cycle's
accepted output, with the cycle's iteration trajectory, method
refinement, per-source coverage analysis, and acceptance criterion
reported separately in {[}4{]}.

\subsubsection{1.2 What this paper
publishes}\label{what-this-paper-publishes}

This paper publishes the v1 artefact and reports three contributions,
each verifiable against the artefact and its validation outputs:

\textbf{Contribution 1 --- AppSec Core v1 as a formalized artefact.} v1
is delivered as an OWL 2 DL ontology export (path under SHA-256 pinning
at the cycle-close release tag) and a SHACL Core constraint set
published as a complementary apparatus (§4). The ontology declares the
same ten slices as v0 (§2), defines class hierarchies for the four
populated entity types (ControlObjective, Practice, Mechanism, Artifact)
plus the EvidencePattern declarative class with its controlled
vocabularies (§3), and is populated with 259 typed instances. The SHACL
apparatus comprises six schema-derived ontology shapes (regenerable from
the canonical YAML schema by a published build script under SHA-256
pinning) and five hand-maintained consumer-conformance shapes encoding
model invariants. SHACL execution against the populated graph yields
\texttt{sh:conforms\ =\ true} with zero \texttt{sh:Violation} under both
the SHACL Core reference-implementation validator \texttt{pyshacl}
(version 0.31.0 with \texttt{rdflib} 7.6.0) and an in-house
bounded-subset validator at parity for the six ontology shapes (§4).

\textbf{Contribution 2 --- Schema preservation under multi-source
pressure.} The ten slices, the four populated entity types, the
EvidencePattern declarative class, and the cross-slice relation model
declared in v0 are preserved unchanged in v1 (§6). Under the corpus
expansion absorbed by the cycle reported in {[}4{]} --- from the five
first-wave sources {[}2{]} to thirty-one sources at cycle close (a
6.2-fold source-count expansion; the cycle's substrate at close records
3,861 source-items decomposed into 18,673 normalized claims) --- the
transition from v0 to v1 absorbed the expansion as \textbf{additive
instance population}: per-entity counts grew (ControlObjective 70→75;
Practice 63→69; Mechanism 48→58; Artifact 53→57; total typed instances
234 → 259, a +25 net delta), but no slice was added, removed, or
repartitioned, no populated entity type was added, the EvidencePattern
declarative class was retained without modification, and no cross-slice
relation was added or modified. The v0→v1 entity-count delta decomposes
by attribution to the cycle's worked governance decisions (§6, §8, §9):
thirteen additions through ACR-001's promotion (+3 CO + 3 P + 3 M + 4
A), five through ACR-002's promotion (+1 CO + 2 P + 2 M), three through
ACR-004's promotion (+1 CO + 1 P + 1 M), with four further within-slice
refinement additions (Mechanism-level only) documented per-instance in
the artefact's release notes. Whether the schema continues to absorb
additional expansion additively under future source pressure is an
empirical hypothesis carried by the cycle's reusable workflow {[}4{]};
this paper does not claim future-source bounding, only that the schema
absorbed the cycle-close corpus as additive population.

\textbf{The substantive engineering hypothesis the artefact embodies
(forward-looking).} C2's evidence supports a stronger hypothesis the
programme carries --- that any new normative framework F (whether
technical standard, regulatory regime, or industry guidance) decomposes
into (a) claims mapping to existing v1 typed entities, normalising into
actionable engineering substance, and (b) residual content routed to the
Manual's prose surface, without ontology schema disruption, modulo
ACR-warranted additive amendments under the protocol of §7. This is not
claimed as an established result of this paper. The framework-horizontal
stress test --- explicit decomposition against current and emerging
regulatory frameworks (DORA, NIS2, CRA, GDPR, PCI-DSS, HIPAA, NIST SP
800-53, and analogous regimes) --- is registered as future work for the
programme; v1 provides the typed normalisation target the test will
exercise, with {[}4{]}'s reusable workflow as the methodological
discipline under which each absorption is governed.

\textbf{Contribution 3 --- The AppSec Core Change Request (ACR) protocol
with its explicit four-condition promotion threshold, demonstrated by
four symmetric worked decisions.} The ACR protocol is the governance
instrument by which candidate changes to the bounded ontology are
admitted or refused. It is published in this paper because it governs
the v0→v1 entity-count delta documented under C2 and because the four
worked decisions are part of v1's published content. The protocol's
threshold (§7) is stated in four conditions: (i) \textbf{multi-source
convergence} --- the candidate's substance appears across at least five
independent sources drawn from at least three distinct organisational
authorities; (ii) \textbf{multi-method convergence} --- the candidate is
detected by independent normalization signals during the cycle's
re-application {[}4{]}; (iii) \textbf{practitioner-Manual content
backing} --- the candidate's substance has a counterpart in the
prescriptive practitioner Manual that complements the ontology; (iv)
\textbf{slice fit} --- the candidate fits within an existing slice's
contract or under a defensible additive amendment to that contract.
Symmetric application across four cases --- three promotions (ACR-001
\emph{Secure Configuration Baseline Integrity} in §8.1; ACR-002
\emph{Security Requirements Lifecycle Management} in §8.2; ACR-004
\emph{Output Rendering Safety / Context-Aware Encoding} in §9.2) and one
non-admission (ACR-003 \emph{Multi-anchor adjacency candidate batch} in
§9.1) --- establishes that the protocol's threshold is testable rather
than retrofitted: the same conditions produce non-promotion under
measurement just as they produce promotion. The non-admission consumes
its identifier as a discipline against governance-history compression.

These three are \emph{factual} claims demonstrable by inspection of the
published artefact, the published validation reports, and the four
worked decisions. The paper does \textbf{not} absorb the design science
cycle's contribution as its own (the iteration trajectory, the method
refinement from keyword-first to multistage, and the per-source coverage
analysis are reported in {[}4{]}) and does \textbf{not} absorb the
Manual-validation cycle's contribution (the bidirectional validation of
v1's typed substrate against the practitioner Manual's prose layer is
reported in {[}5{]}).

\subsubsection{1.3 The two surfaces of the programme architecture
(context)}\label{the-two-surfaces-of-the-programme-architecture-context}

The artefact published here lives within the \textbf{Security-by-Design
Theory of Everything (SbD-ToE)} programme architecture {[}3{]}, which
the reader needs to understand only at the level of background context
for this paper's claims. The SbD-ToE programme delivers
application-security guidance through two complementary surfaces with
asymmetric scope. The \textbf{bounded ontology} (this paper's artefact)
is a typed, validated, machine-readable subset of the
application-security space, bounded by per-slice contracts and by the
ACR governance protocol of §7, addressing the
development-and-maintenance substance an engineer must satisfy at
engagement scale (per pull request, per sprint, per release). The
\textbf{practitioner Manual} is comprehensive: it carries both the
engagement-view substance --- a subset of which the ontology normalizes
into typed entities, with the rest covered in prose throughout the
Manual --- and the organizational-view substance such as governance,
training, vendor contractual security management, regulatory compliance
regimes, and workforce programmes that lies outside the ontology's
bounds by design. The ontology is, in this sense, a normalized subset of
part of what the Manual covers; the Manual is the comprehensive
prescriptive corpus that complements the ontology in the programme
architecture. This paper publishes the ontology surface; the
bidirectional Manual-validation cycle is the subject of {[}5{]}. The
downstream apparatus papers {[}6, 7, 8{]} consume both surfaces in
combination as a security-requirements harness for code-generation
consumers.

This paper does not need the reader to know more than the above about
the Manual surface; the schema-preservation claim of C2 and the ACR
protocol of C3 are stated at the ontology's level. The third condition
of the ACR threshold (practitioner-Manual content backing) references
the Manual as a content-correctness checkpoint, not as a substantive
contribution of this paper.

\textbf{Engineering substance vs governance substance --- the
design-intentional split.} The bounded ontology absorbs the
\textbf{engineering substance} of normative AppSec sources (typically:
secure development practices, technical safeguards, vulnerability
management, incident detection, supply-chain integrity, configuration
baseline integrity, identity and access controls, input/output
validation, output encoding, threat modelling, security-event logging).
The \textbf{governance substance} that frequently accompanies normative
sources --- board-level reporting obligations, vendor contractual risk
management, regulatory examination response, workforce training
programmes, organisation-level policy authorities, audit-firm interface
--- is routed by design to the Manual's prose surface, not absorbed into
the ontology's typed entities. This split is intentional, not
incidental: the ontology's purpose is to provide a typed normalisation
target for engineering substance an engineer (or LLM/MCP-grounded coding
agent) can act upon at the application's development scale; the
governance substance complementing that engineering work has a different
operational consumer (programme management, governance committees,
organisational compliance functions) and lives in the Manual's prose
accordingly. The substantive claim of the programme is therefore not
``v1 substitutes reading the framework'' --- it is ``v1 provides the
normalisation target for the framework's engineering substance, the
Manual complements for the framework's governance substance, and the ACR
protocol governs where engineering substance pressures schema.''

\textbf{Illustrative consumption shape (engagement scale).} A
representative consumption pattern: an engineer or coding agent receives
a code change touching authentication middleware. Querying v1 with the
change's substantive surface returns a typed working set --- for
example, ControlObjectives in the Identity, Access \& Session Trust
slice (\texttt{ACO-IAT-NNN} family); their associated Practices
(\texttt{ACP-IAT-NNN}); the supporting Mechanisms that operationalise
those Practices (\texttt{ACM-IAT-NNN}); and the expected Artifacts
(\texttt{ACA-IAT-NNN}) the change should produce or update for evidence.
The engineer's task collapses from ``read the relevant chapters of every
authority'' to ``demonstrate satisfaction of this typed working set
against the diff under review'', with the Manual's prose surface
remaining available for the governance substance that does not normalise
into typed entities. This consumption pattern is the
\emph{engagement-scale operationalisation} the ontology's typing
supports; the empirical evaluation of this consumption pattern under
retrieval, language-model grounding, and grounded-versus-ungrounded
compilation is reported in {[}6, 7, 8{]} respectively. This paper's
contribution is the publication of v1 as the typed reference layer; the
operational evaluation lives downstream.

\subsubsection{1.4 Scope}\label{scope}

This paper publishes the v1 artefact and reports the three contributions
above. The following are out of scope and are reported separately:

\begin{itemize}
\tightlist
\item
  \emph{The design science cycle that produced v1 from v0} --- its
  iteration trajectory, the method refinement (keyword-first →
  multistage normalization), the cross-source pressure analysis, the
  cycle's \emph{good-for-intended-fit} acceptance criterion and its
  evidence pillars, and the per-source-into-Reference mapping skin
  compatible with the NIST OLIR programme templates {[}32, 33{]} --- is
  reported in {[}4{]}.
\item
  \emph{Per-source corpus analysis} of the thirty-one sources from which
  v1's instance population was drawn, including per-source convergence
  rates, per-source adjacency distributions, and override-mechanism logs
  --- is reported in {[}4{]}.
\item
  \emph{The bidirectional validation of v1's typed substrate against the
  practitioner Manual's prose layer} --- including the Manual-coverage
  cycle, the edit-versus-traceback decisioning for each ACR addition,
  and the closure event that freezes corpus + Manual + ontology jointly
  for downstream consumption --- is reported in {[}5{]}.
\item
  \emph{Retrieval, large-language-model grounding, and experimental
  evaluation} against v1 --- are addressed by the programme's papers
  {[}6{]}, {[}7{]}, and {[}8{]}.
\end{itemize}

The empirical case for v1 in this paper is restricted to what the
artefact and its validation directly demonstrate: structural conformance
(§4), schema preservation across the v0→v1 transition (§6), and the
four-decision symmetric application of the ACR protocol (§7--§9). Claims
of correctness for individual instance mappings, peer review of the
per-source corpus, and stakeholder evaluation of v1's downstream utility
are open and registered as future work (§10).

\textbf{Bounded engineering-ontology posture.} AppSec Core v1
intentionally adopts a bounded engineering-ontology posture rather than
attempting deep ontological realism, upper-ontology completeness, or
richly inferential semantic-web research. The artefact's purpose is to
support operational normalization of heterogeneous AppSec sources at
engineering scale (the harness consumer of {[}3, 6{]}), not to construct
a foundational semantic theory of cybersecurity. In the terminology of
Guarino {[}23{]}, AppSec Core is an \textbf{application ontology} ---
domain- and task-specific, derivation-grounded rather than
foundation-grounded, with no commitment to upper-ontology categorial
alignment (BFO, DOLCE, UFO). The OWL 2 DL ontology provides typed
entities, an explicit cross-slice relation type system, controlled
vocabularies, and the per-slice contracts that bound each slice's scope
and acceptance conditions; richer inferential commitments (deep
subsumption hierarchies, role compositions beyond the declared
properties, non-trivial logical entailment over the populated graph) are
out of scope by design. The artefact is positioned as a \textbf{bounded
normalization ontology} rather than as a candidate for semantic-web
representational completeness.

\textbf{Why an ontology rather than a typed knowledge graph.} The choice
of OWL 2 DL with SHACL apparatus rests on four substantive architectural
requirements. First, the bounded ontology operates as a \textbf{typed
reference for heterogeneous-source normalization}: each external
source's items are mapped against typed entities (ControlObjective /
Practice / Mechanism / Artifact) under explicit per-slice contracts, and
the OWL class hierarchy + property domain/range declarations make the
typing computationally enforceable rather than convention-only. Second,
the ACR governance protocol of §7 operates over \textbf{declared
structural commitments} (slice partition, populated entity types,
cross-slice relation type set, per-slice contracts); without OWL-level
commitments, ``schema preservation'' would be a property of editorial
discipline rather than of the artefact's formal surface. Third, the
SHACL apparatus of §4 enforces the cross-instance model invariants at a
\textbf{closed-world constraint layer} complementing OWL's open-world
class semantics; this dual layering is operationally valuable for
downstream tooling (validators, retrieval consumers) and is not
satisfied by typed-KG conventions that lack a constraint-checking layer.
Fourth, the bounded ontology is consumed by downstream MCP-grounded
coding agents {[}8{]} and retrieval consumers {[}6, 7{]} requiring
\textbf{RDF/SPARQL-native serialisation with closed-world constraint
enforcement at query time}; this consumption profile is satisfied
natively by OWL 2 DL + SHACL Core. Lighter-weight alternatives were
considered --- typed knowledge graphs (RDF + property-graph conventions
without OWL semantics or SHACL constraint enforcement) and schema-only
modelling layers (LinkML; JSON Schema bound to SHACL post-hoc) --- but
the four requirements above are not satisfied natively by these
alternatives and would require non-trivial tooling supplementation; OWL
2 DL + SHACL Core was therefore considered the most appropriate form for
the artefact's normalization purpose, governance commitments,
constraint-layer enforcement, and downstream consumption profile.

\textbf{Positioning of v1.} This paper does not claim that AppSec Core
v1 is a complete or semantically final ontology of application security;
it publishes a bounded, formally constrained, governance-controlled
normalization ontology whose v0→v1 evolution absorbed the cycle-close
corpus without schema restructuring.

\subsubsection{1.5 Related ontology work and methodology
positioning}\label{related-ontology-work-and-methodology-positioning}

AppSec Core v1 occupies a complementary layer to existing cybersecurity
ontologies. The Unified Cybersecurity Ontology (UCO) {[}24{]} and the
STUCCO ontology {[}25{]} integrate threat-intelligence schemas (CVE,
CWE, CAPEC, STIX, and analogous structured-threat surfaces) for
situational-awareness, threat-correlation, and intelligence-fusion use
cases; recent cybersecurity knowledge-graph work continues to extend
that surface {[}26{]}. AppSec Core targets a different operational
consumer: the engineering-practice substance of heterogeneous AppSec and
regulatory sources, normalized for development-and-maintenance
compliance reuse at engagement scale (§1.3). The two ontology families
address complementary layers of the cybersecurity-knowledge stack ---
threat intelligence (UCO/STUCCO/CKG) and engineering-practice
normalization (AppSec Core) --- and can be combined in downstream
tooling without semantic conflict.

The cycle that produced v1 {[}4{]} instantiates an iterative ontology
engineering workflow conceptually aligned with the NeOn methodology
{[}27{]} (scenario-based development with formal change control) and the
LOT methodology {[}28{]} (industrial-oriented engineering with explicit
governance instruments), with the ACR governance protocol of §7
providing the specific change-control instrument adapted to bounded
normalization ontologies. The programme's design science discipline
follows the canonical DSR framework of Hevner et al.~{[}20{]} and
Wieringa {[}21{]} under the methodological reporting standards of Ralph
et al.~{[}22{]}, with Peffers et al.~{[}29{]} providing the
operationalisation pattern most often cited in information-systems DSR.

\subsubsection{1.6 Document structure}\label{document-structure}

The rest of the paper is organised as follows. §2 describes the v1 slice
and instance structure. §3 specifies the OWL 2 formalization. §4
specifies the SHACL apparatus and reports the validation outcome. §5
documents v1 as the accepted output of the cycle reported in {[}4{]}. §6
reports the schema-preservation evidence across the v0→v1 transition
(Contribution 2). §7 specifies the ACR protocol and its four-condition
promotion threshold (Contribution 3). §8 reports the worked decisions on
candidates carried forward from the v0 baseline (ACR-001 and ACR-002,
both promoted). §9 reports the worked decisions on candidates surfaced
by the cycle's re-application of the refined method (ACR-003, not
admitted; ACR-004, promoted). §10 states limitations. §11 specifies the
reproducibility surface. §12 lists references.

\begin{center}\rule{0.5\linewidth}{0.5pt}\end{center}

\subsection{2. AppSec Core v1: structure and instance
population}\label{appsec-core-v1-structure-and-instance-population}

\subsubsection{2.1 The ten-slice
partition}\label{the-ten-slice-partition}

AppSec Core v1 preserves the ten-slice partition introduced in v0 {[}1,
§3{]}. v0's partition was derived bottom-up from the practitioner Manual
corpus through expert qualitative reading {[}1, §3.2{]}, with each slice
opened serially under an explicit per-slice contract specifying its
scope, its non-goals, and its boundary conditions {[}1, §3.5{]}. Each
slice has a declared primary control surface --- the operational
governance locus at which the slice's controls converge in the delivered
software lifecycle --- captured by a categorical descriptor
(\texttt{single\_broad\_control\_surface},
\texttt{dual\_control\_surface},
\texttt{multi\_control\_convergent\_surface},
\texttt{split\_control\_surface}).

The slice partition is presented in Table 1 with the v1 instance
population per slice. Slice scope summaries are abridged from the
per-slice contracts pinned at the cycle-close release tag; the contracts
are the canonical reference for the full scope, non-goals, and boundary
conditions of each slice.

\textbf{Table 1.} AppSec Core v1 slice registry, with scope summary and
v1 typed-instance population. Slice scope summaries are abridged from
each slice's per-slice contract at the cycle-close release tag
(\texttt{ontology-v1-final}); the contracts are the canonical reference
for the full scope, non-goals, and boundary conditions of each slice.
Slice abbreviation is the entity-prefix convention used throughout the
ontology's identifier scheme (e.g., \texttt{ACO-RPR-NNN} =
ControlObjective in the RPR slice). Population per slice is the count of
ControlObjectives + Practices + Mechanisms + Artifacts populated under
the slice's per-slice contract at the cycle-close release tag, read
directly from \texttt{ontology-v1-final}'s instance index. Slice IDs
follow the v0-published convention {[}1{]}.

\begin{longtable}[]{@{}
  >{\raggedright\arraybackslash}p{(\linewidth - 6\tabcolsep) * \real{0.2500}}
  >{\raggedright\arraybackslash}p{(\linewidth - 6\tabcolsep) * \real{0.2500}}
  >{\raggedright\arraybackslash}p{(\linewidth - 6\tabcolsep) * \real{0.2500}}
  >{\raggedright\arraybackslash}p{(\linewidth - 6\tabcolsep) * \real{0.2500}}@{}}
\toprule\noalign{}
\begin{minipage}[b]{\linewidth}\raggedright
Slice ID
\end{minipage} & \begin{minipage}[b]{\linewidth}\raggedright
Slice abbreviation
\end{minipage} & \begin{minipage}[b]{\linewidth}\raggedright
Scope (abridged)
\end{minipage} & \begin{minipage}[b]{\linewidth}\raggedright
v1 population
\end{minipage} \\
\midrule\noalign{}
\endhead
\bottomrule\noalign{}
\endlastfoot
ASC-01 & SCBI & Supply Chain \& Build Integrity (build pipeline
integrity, source repository hygiene, dependency hygiene,
container/image supply-chain controls) & 30 (7 CO + 7 P + 6 M + 10 A) \\
ASC-02 & IAT & Identity, Access \& Session Trust (authentication,
authorisation, session and identity boundary controls) & 25 (7 + 6 + 6 +
6) \\
ASC-03 & ATB & Architecture \& Trust Boundaries (architectural trust
boundary identification, sandboxing, isolation) & 23 (7 + 7 + 5 + 4) \\
ASC-04 & TSV & Testing, Security Validation \& Empirical Assurance
(SAST, DAST, penetration testing, security-relevant unit testing) & 25
(7 + 7 + 5 + 6) \\
ASC-05 & TMR & Threat Modeling, Risk Disposition \& Mitigation
Traceability (threat identification, risk disposition, mitigation
traceability; ACR-002 additive amendment adds
requirements-as-living-artefacts dimension per §8.2) & 29 (8 + 9 + 8 +
4; includes ACR-002 additions) \\
ASC-06 & SPC & Secret Handling, Protected Configuration \& Operational
Identities (secret material handling, protected configuration values,
operational-identity discipline) & 19 (7 + 6 + 4 + 2) \\
ASC-07 & IVF & Input Validation, Safe Parsing \& Controlled Failure
(input boundary discipline, safe parsing, controlled failure under
hostile or malformed input; ACR-004 additive amendment expands scope to
input/output validation and controlled failure per §9.2) & 26 (8 + 7 + 5
+ 6; includes ACR-004 additions) \\
ASC-08 & ITS & Integration Trust \& Service-to-Service Security
(service-to-service authentication, integration boundary controls, trust
delegation) & 23 (7 + 6 + 5 + 5) \\
ASC-09 & RPR & Release Promotion, Controlled Rollout \& Rollback
Readiness (release-promotion control, secure configuration baseline
integrity, controlled rollout and rollback discipline; ACR-001 additive
amendment adds baseline-integrity ControlObjectives per §8.1) & 39 (10 +
9 + 8 + 12; includes ACR-001 additions) \\
ASC-10 & SLG & Security Event Logging, Audit Trail \& Centralized
Logging (audit-log integrity, security-event recording, log retention
and review) & 20 (7 + 5 + 6 + 2) \\
\end{longtable}

Per-slice population counts are reported in the cycle-close release
tag's slice registry artefact. The populated graph at v1 contains 259
typed instances total: 75 ControlObjectives + 69 Practices + 58
Mechanisms + 57 Artifacts.

\subsubsection{2.2 The four populated entity types and the declarative
EvidencePattern
class}\label{the-four-populated-entity-types-and-the-declarative-evidencepattern-class}

v1 preserves the four populated entity types declared in v0:

\begin{itemize}
\tightlist
\item
  \textbf{ControlObjective}: a reusable domain constraint or assurance
  goal, formulated in the verification voice (``the system
  shall\ldots{}'');
\item
  \textbf{Practice}: a reusable human or process discipline through
  which control objectives are realized;
\item
  \textbf{Mechanism}: a concrete, enforceable means by which a practice
  is implemented (a tool, a configuration, a runtime control);
\item
  \textbf{Artifact}: a control-relevant or evidence-bearing output (a
  document, a configuration record, a verifiable signed artefact).
\end{itemize}

\textbf{Identifier convention.} Every populated entity carries an
identifier of the form
\texttt{AC\textless{}T\textgreater{}-\textless{}SLICE\textgreater{}-\textless{}NNN\textgreater{}},
where:

\begin{itemize}
\tightlist
\item
  \texttt{AC} is the AppSec Core programme prefix;
\item
  \texttt{\textless{}T\textgreater{}} is the entity-type initial:
  \texttt{O} for \textbf{C}ontrol\textbf{O}bjective, \texttt{P} for
  \textbf{P}ractice, \texttt{M} for \textbf{M}echanism, \texttt{A} for
  \textbf{A}rtifact;
\item
  \texttt{\textless{}SLICE\textgreater{}} is the slice abbreviation per
  Table 1 (\texttt{SCBI}, \texttt{IAT}, \texttt{ATB}, \texttt{TSV},
  \texttt{TMR}, \texttt{SPC}, \texttt{IVF}, \texttt{ITS}, \texttt{RPR},
  \texttt{SLG});
\item
  \texttt{\textless{}NNN\textgreater{}} is a sequential three-digit
  per-type-per-slice index.
\end{itemize}

Examples: \texttt{ACO-IAT-001} is the first ControlObjective in the
Identity, Access \& Session Trust slice (ASC-02); \texttt{ACP-RPR-008}
is the eighth Practice in the Release Promotion, Controlled Rollout \&
Rollback Readiness slice (ASC-09); \texttt{ACM-IVF-005} is the fifth
Mechanism in the Input Validation, Safe Parsing \& Controlled Failure
slice (ASC-07); \texttt{ACA-SCBI-001} is the first Artifact in the
Supply Chain \& Build Integrity slice (ASC-01). Slice instances
themselves use identifier \texttt{Slice\textless{}ASC-NN\textgreater{}}
aligned with the slice IDs of Table 1.

\textbf{Canonical source.} The identifier scheme is \textbf{declared
canonically} in the YAML schema source at the cycle-close release tag
(\texttt{appsec-core-entity-schema-v0-draft.yaml} within the bundle's
\texttt{ontology-v1-final} artefact, §11.1); the OWL TTL export
(\texttt{appsec-core-v0-bounded-v1.ttl}, also within
\texttt{ontology-v1-final}) \textbf{instantiates} the scheme as full
IRIs under the \texttt{ac:} namespace, using the underscore form
\texttt{AC\textless{}T\textgreater{}\_\textless{}SLICE\textgreater{}\_\textless{}NNN\textgreater{}}
per TTL serialisation conventions (e.g.,
\texttt{ac:ACO\_RPR\_008\ a\ ac:ControlObjective}). The body prose of
this paper uses the hyphenated form
\texttt{AC\textless{}T\textgreater{}-\textless{}SLICE\textgreater{}-\textless{}NNN\textgreater{}}
for readability; the two are semantically equivalent and
identifier-stable across the bundle's serialisation formats.

A fifth class, \textbf{EvidencePattern}, was introduced in v0 as a
supporting indexed class expressing the expected evidence shape for
deterministic review of artefact compliance. v0 published a supporting
index of 213 evidence-pattern entries {[}1, §6{]}. \textbf{v1's
populated graph contains zero EvidencePattern instances.} What v1
preserves is the EvidencePattern \emph{class declaration} --- its formal
class hierarchy, OWL property declarations, and SHACL
\texttt{EvidencePatternShape} (with
\texttt{sh:targetClass\ ac:EvidencePattern} plus optional property
shapes referencing it from ControlObjective and Artifact) --- at the
apparatus level; the v0 supporting index of 213 patterns is preserved as
a \emph{separate artefact} within the cycle's deposited bundle
(\texttt{appsec-core-evidence-pattern-index-v0.yaml}, see §11) and is
consumable by downstream tooling that requires evidence-pattern
semantics, but it is not migrated into v1's populated graph as
first-class typed instances. The class is intentionally unpopulated at
v1: the SHACL conformance report records
\texttt{target\_node\_count\ =\ 0} for \texttt{EvidencePatternShape}.
Population of the class as first-class typed instances at the populated
graph is registered as future programme work (§10.4). The 234 / 259
figures reported throughout this paper for the v0 / v1 typed-instance
counts therefore refer to the four populated entity types only and
exclude the EvidencePattern declarative class.

\subsubsection{2.3 The cross-slice relation
model}\label{the-cross-slice-relation-model}

v1 preserves the cross-slice relation model of v0 {[}1, §3.6{]}. The OWL
ontology declares the relations under the canonical names exported by
the build pipeline at \texttt{ontology-v1-final}: every populated
ControlObjective is linked to at least one Practice
(\texttt{objective\_realized\_by\_practice}) and at least one Mechanism
(\texttt{objective\_implemented\_by\_mechanism}); ControlObjectives
expect supporting Artifacts (\texttt{objective\_expects\_artifact});
slice membership is captured by \texttt{belongsToSlice}. The relation
TYPES (objective↔practice, objective↔mechanism, objective↔artifact,
slice membership) and connections between entity types are preserved
unchanged from v0; the relation NAMES exported by the v1 OWL build
pipeline use mixed convention (\texttt{belongsToSlice} camelCase
short-form for slice membership; \texttt{objective\_*\_by\_*} and
\texttt{objective\_expects\_*} snake\_case long-form for the typed
inter-entity relations) consistent with the v1 OWL ontology release. No
new cross-slice relation type has been added at v1.

\textbf{Slice-membership functionality and cross-slice substantive
coupling.} \texttt{belongsToSlice} is \textbf{functional}: each
populated entity (ControlObjective / Practice / Mechanism / Artifact)
belongs to exactly one slice, expressed at the SHACL apparatus level as
\texttt{sh:maxCount\ 1} on the property's shape. Substantive content
that interacts across slices --- for example, output rendering safety
(ACR-004's substance) plausibly touches the IVF slice (input/output
validation), the SCBI slice (build-time encoding configuration), and the
ITS slice (service-to-service trust at the integration boundary) --- is
modelled via cross-slice \textbf{relations} (objective↔practice,
objective↔mechanism, objective↔artifact triples that may cross slice
boundaries), \textbf{not} via multi-slice membership. The canonical
placement discipline is: each entity is placed in the slice of its
\textbf{primary control surface} (the operational governance locus at
which the entity converges in the delivered software lifecycle), with
cross-slice substantive coupling expressed through the typed relations
that explicitly connect entities across slice boundaries. This
discipline avoids role conflation (the same entity playing structurally
different roles across slices) by reserving the primary placement for
the slice where the entity's primary control surface lives, and using
the cross-slice relations as the operational language for inter-slice
coupling. ACR-004's IVF placement instantiates this discipline: the
substance's primary control surface is the input/output validation
boundary at runtime; coupling to SCBI (build-time encoding
configuration) and to ITS (service-to-service rendering boundaries) is
expressed through \texttt{objective\_implemented\_by\_mechanism} triples
whose Mechanism instances may carry their own primary placement in those
adjacent slices.

\begin{center}\rule{0.5\linewidth}{0.5pt}\end{center}

\subsection{3. OWL 2 formalization}\label{owl-2-formalization}

\subsubsection{3.1 Class hierarchy}\label{class-hierarchy}

The OWL 2 DL ontology declares classes for the four populated entity
types (ControlObjective, Practice, Mechanism, Artifact) and for the
declarative EvidencePattern class, plus a Slice class whose individuals
are the ten populated slices. The class hierarchy is shallow: each typed
entity is a direct subclass of the top-level type, with no subtyping by
slice (slice membership is captured by an object property, not by class
hierarchy). The Slice class carries a \texttt{current\_control\_mode}
data property whose values are restricted to the four
controlled-vocabulary terms (\texttt{single\_broad\_control\_surface},
\texttt{dual\_control\_surface},
\texttt{multi\_control\_convergent\_surface},
\texttt{split\_control\_surface}). The ontology lies within the OWL 2 DL
profile; classification and consistency checking are decidable in
NEXPTIME under standard reasoner implementations (HermiT, Pellet,
OpenLLET).

\textbf{OWL 2 profile choice.} The DL profile was selected over EL, QL,
and RL on three grounds (note: all three OWL 2 sub-profiles support
\texttt{rdfs:domain} and \texttt{rdfs:range} axioms; the differences
across profiles concern axiom families that combine class and property
expressivity in richer ways). (a) The artefact uses class disjointness
and complement-style axioms for the populated entity-type partition
(ControlObjective / Practice / Mechanism / Artifact / EvidencePattern as
mutually disjoint classes) --- EL excludes negation and disjointness; QL
restricts class axiom forms in ways that would force lossy translation
of the disjointness structure. (b) The artefact's reasoning use case is
closed-world structural conformance checking at the SHACL apparatus
layer (§4) rather than open-world query answering against very large
ABoxes (which would favour QL); the OWL-level reasoning uses are
restricted to class-hierarchy classification and consistency checking,
well within DL expressivity. (c) The artefact is consumed by downstream
tooling that combines OWL-level class-axiom reasoning with SHACL-level
instance-graph constraint checking; DL provides the minimum expressivity
that supports the disjointness axioms above without forcing translation
to a profile whose restrictions are not motivated by the artefact's
reasoning use cases.

\textbf{Class-level disjointness at the vocabulary-value / entity
boundary.} Beyond the entity-type partition disjointness that motivates
the DL profile choice above, v1 declares a class-level
\texttt{owl:disjointWith} axiom between
\texttt{ac:ControlledVocabularyValue} (the abstract grouping superclass
for the six controlled-vocabulary classes --- \texttt{ac:ObjectiveKind},
\texttt{ac:ObjectiveType}, \texttt{ac:ArtifactRole},
\texttt{ac:PracticeFamily}, \texttt{ac:MechanismFamily},
\texttt{ac:DetectableSurface}) and \texttt{ac:AppSecCoreEntity} (the
abstract grouping superclass for the four populated entity types plus
the EvidencePattern declarative class). The axiom formalises the
categorical separation between vocabulary-value classes and AppSec Core
entity classes at the OWL semantics level. It preserves OWL 2 DL
decidability and is reachable to reasoners during class-hierarchy
verification.

\subsubsection{3.2 Object and data
properties}\label{object-and-data-properties}

The ontology declares object properties for the cross-slice relations
(\texttt{belongsToSlice}, \texttt{objective\_realized\_by\_practice},
\texttt{objective\_implemented\_by\_mechanism},
\texttt{objective\_expects\_artifact}, plus the EvidencePattern indexing
properties declared but not exercised at populated-graph level under
v1). Data properties capture the per-instance descriptive content (item
identifier, slice-scoped name, descriptive narrative,
controlled-vocabulary tags). All properties are declared with explicit
domain and range restrictions; cardinality constraints are expressed at
the SHACL apparatus level (§4) rather than as OWL cardinality axioms, to
keep the OWL ontology decidable under DL while expressing the
populated-graph constraints in the SHACL layer where they are
operationally consumed.

\subsubsection{3.3 Controlled
vocabularies}\label{controlled-vocabularies}

Controlled vocabularies cover (a) slice control modes (the four-term
enumeration above), (b) per-instance landing-strength labels
(\texttt{direct}, \texttt{strong}, \texttt{bounded}) used by downstream
consumers to communicate the cycle's mapping confidence at instance
level, and (c) cross-slice vocabulary terms (a curated set of
approximately fifty domain-vocabulary terms covering the recurrent
technical-anchor vocabulary across the AppSec slices). The controlled
vocabularies are declared as OWL annotation properties with the term
values pinned in the ontology; SHACL shapes (§4) enforce closed-world
membership where required.

\subsubsection{3.4 Instance population}\label{instance-population}

At the populated graph for v1, the OWL ontology carries 259 typed
instances across the four populated entity types as enumerated in §2.1.
The bounded Turtle export (\texttt{ontology-v1-final}) serialises
\textbf{1,970 data triples} at the cycle-close release tag. Each
instance has a unique IRI under the cycle's release tag namespace;
relations between instances are explicit triples in the populated graph.
Reasoner-decidable class hierarchy and property-domain/range checks pass
under standard OWL 2 DL reasoners; the populated graph is consumable by
reasoners that implement the OWL 2 DL profile.

\begin{center}\rule{0.5\linewidth}{0.5pt}\end{center}

\subsection{4. SHACL formalization (the
apparatus)}\label{shacl-formalization-the-apparatus}

\subsubsection{4.1 Apparatus composition}\label{apparatus-composition}

The SHACL apparatus accompanying v1 is composed of two complementary
shape sets, an architectural pattern published as part of the cycle's
release {[}4{]}:

\begin{itemize}
\tightlist
\item
  \textbf{Schema-derived ontology shapes} (file:
  \texttt{appsec-core-v0-shapes.ttl} within the bundle's apparatus
  artefact). Six NodeShapes covering the five typed entity classes
  (Slice, ControlObjective, Practice, Mechanism, Artifact) plus the
  EvidencePattern declarative class. These shapes are
  \textbf{regenerable} from the canonical YAML schema by a published
  build script (\texttt{build\_shacl.py} at the same release tag);
  regenerating from the YAML under the documented build environment
  produces a deterministic SHACL output reproducible against the
  released shapes file (isomorphic up to blank node labeling under the
  RDFLib serialisation defaults of the documented environment).
\item
  \textbf{Consumer-conformance shapes} (file:
  \texttt{consumer-conformance-shapes.ttl} within the bundle's apparatus
  artefact). Five \texttt{ac:Claim}-typed shapes encoding model
  invariants that the populated graph must satisfy beyond the
  schema-derived shapes' per-entity field-level constraints. These five
  shapes are \textbf{hand-maintained} from the model-invariant decisions
  recorded with the cycle {[}4{]} and capture: (M1') slice-coherence
  invariants for each populated slice; (M3) consistency between Practice
  and ControlObjective populations within each slice; (M4) consistency
  between Mechanism and ControlObjective populations through the chain
  Mechanism → Practice → ControlObjective; (M4-card) claim
  well-formedness invariants; and (M4-card) claim target
  referential-integrity invariants. The hand-maintained shapes
  complement the schema-derived shapes at the populated-graph level: the
  schema-derived shapes verify that each instance carries its required
  fields under its type-level contract, and the hand-maintained shapes
  verify that the relations across instances satisfy the model
  invariants.
\end{itemize}

The apparatus composition is published under a single composed shapes
graph at validation time; the two shape files compose without conflict.

\subsubsection{4.2 Validation under two complementary
validators}\label{validation-under-two-complementary-validators}

The populated graph is validated against the composed shapes graph under
two complementary validators:

\begin{itemize}
\tightlist
\item
  The \textbf{SHACL Core reference-implementation validator}
  \texttt{pyshacl} at version 0.31.0 with \texttt{rdflib} 7.6.0.
  \texttt{pyshacl} implements the SHACL Core specification and is used
  here as the canonical validation oracle for the populated graph; the
  implementation is community-maintained against the W3C SHACL
  Recommendation rather than W3C-endorsed.
\item
  An \textbf{in-house bounded-subset validator} at parity with
  \texttt{pyshacl} for the apparatus's six ontology shapes. The
  bounded-subset validator was developed as a programme-internal tool
  predating the apparatus's \texttt{pyshacl}-canonical release and has
  been retained as a parallel validation surface for tooling that
  consumes per-shape granular validation outputs (the bounded-subset
  validator emits per-shape conformance signals usable by downstream
  tooling without parsing the full SHACL conformance report). Parity
  with \texttt{pyshacl} is verified at apparatus release time across the
  six ontology shapes and is preserved within the bundle's apparatus
  artefact.
\end{itemize}

The composed SHACL shapes graph carries \textbf{396 shape triples across
11 NodeShapes} (the six schema-derived ontology shapes plus the five
hand-maintained consumer-conformance shapes) at the cycle-close release
tag. Validation against the populated graph at v1's cycle-close release
tag yields:

\begin{itemize}
\tightlist
\item
  \texttt{pyshacl} 0.31.0 + \texttt{rdflib} 7.6.0:
  \texttt{sh:conforms\ =\ true}; zero \texttt{sh:Violation} across the
  composed shapes graph.
\item
  Bounded-subset validator: \texttt{conforms\ =\ true}; zero violations
  across the six ontology shapes (parity scope).
\end{itemize}

The two validators agree at v1's cycle-close. Where the two diverge in
scope, \texttt{pyshacl} is the canonical oracle for the populated-graph
claim; the bounded-subset validator's per-shape granular outputs are
auxiliary to consumer tooling.

\subsubsection{4.3 What the validation evidence
supports}\label{what-the-validation-evidence-supports}

The validation outcome supports the structural-conformance subset of
Contribution 1: at the populated graph for v1, the composed SHACL
apparatus reports conformance under both validators. This is a
\textbf{binary structural conformance} claim. It does \emph{not} claim
that any individual mapping decision recorded against an instance is
content-correct against the source authority that motivated the instance
--- content-correctness of individual mappings is the territory of
subject-matter-expert review, registered as future work (§10). The SHACL
conformance verdict supports that the populated graph satisfies the
apparatus's structural constraints; it does not certify the substantive
correctness of the apparatus's constraints themselves or of the
instances' substantive content.

\textbf{Shape-set completeness disclosure.} The apparatus's evidentiary
value depends on the completeness of the shape set against the
invariants the artefact intends to enforce. The six schema-derived
ontology shapes cover field-level structural constraints per entity type
(required fields, controlled-vocabulary value ranges, identifier
conventions, basic cardinality); the five hand-maintained
consumer-conformance shapes encode the cross-instance model invariants
identified during the cycle's design {[}4{]}. Completeness of the shape
set against \emph{all} possible invariants of a well-formed AppSec Core
ontology is not claimed: specifically, semantic-level invariants (e.g.,
whether a promoted ControlObjective's verification voice is well-formed;
whether a Practice's description is implementable; whether a Mechanism's
enforcement is operationally testable) are not expressible in SHACL Core
and are not captured by the apparatus. The conformance verdict therefore
supports structural well-formedness against the explicit shape set as
released, not semantic correctness or comprehensiveness of the invariant
set itself; expansion of the invariant set to capture additional
semantic constraints is registered as future work (§10).

\begin{center}\rule{0.5\linewidth}{0.5pt}\end{center}

\subsection{5. The accepted output of the design science
cycle}\label{the-accepted-output-of-the-design-science-cycle}

v1 is published here as the \textbf{accepted output} of the design
science cycle reported in {[}4{]}. The cycle reported in {[}4{]}
iterates the v0 ontology and a v0-era keyword-first compilation method
{[}2{]} across three iterations under the pressure of expansion to
thirty-one sources, refining method (Iteration 1), exercising the ACR
protocol of §7 (Iteration 2), and pressure-testing the bounded thesis
under an extension to the AI/ML domain (Iteration 3). The cycle's
acceptance criterion --- \emph{good-for-intended-fit}, operationalised
through four evidence pillars covering structural conformance under
SHACL Core (Pillar 1), per-source coverage acceptable under independent
third-party cross-references (Pillar 2 --- oracle agreement), multi-mode
evaluation returning no failure-class signal (Pillar 3), and rejection
of random source-claim assignment with semantic-specialisation direction
under permutation-test null model (Pillar 4 --- k-way structural
concentration) --- is reported in {[}4{]}; the criterion's first pillar
is satisfied by §4's SHACL conformance verdict.

The artefact reported in this paper is the cycle's ontology-side output
at cycle close: the OWL 2 DL ontology of §3, the SHACL apparatus of §4,
and the populated graph at v1's cycle-close release. The artefact is
SHA-256-pinned within the cycle's deposited bundle (§11.0) under the
stable artefact identifiers \texttt{ontology-v1-final} and
\texttt{apparatus-shacl-pyshacl-v3}. The cycle's narrative is not
re-derived in this paper; the artefact is published as the cycle's
accepted output, with citation to {[}4{]} for the cycle that produced
it.

\begin{center}\rule{0.5\linewidth}{0.5pt}\end{center}

\subsection{6. Schema preservation under multi-source
pressure}\label{schema-preservation-under-multi-source-pressure}

\subsubsection{6.1 What is preserved across
v0→v1}\label{what-is-preserved-across-v0v1}

The following elements of the v0 schema are preserved unchanged in v1:

\begin{itemize}
\tightlist
\item
  The \textbf{ten-slice partition} declared in v0 {[}1, §3{]}. No slice
  was added, removed, or repartitioned at v1; the
  \texttt{current\_control\_mode} of each slice is preserved.
\item
  The \textbf{four populated entity types} (ControlObjective, Practice,
  Mechanism, Artifact). No populated entity type was added.
\item
  The \textbf{EvidencePattern declarative class}. The class declaration
  is preserved at the OWL level; v1 does not populate additional
  EvidencePattern instances at this release.
\item
  The \textbf{cross-slice relation model}: \texttt{belongsToSlice},
  \texttt{objective\_realized\_by\_practice},
  \texttt{objective\_implemented\_by\_mechanism},
  \texttt{objective\_expects\_artifact} (canonical names exported by the
  v1 OWL build pipeline at \texttt{ontology-v1-final}; relation TYPES
  and connections between entity types are preserved unchanged from v0,
  see §2.3). No new cross-slice relation type was added or modified.
\item
  The \textbf{per-slice contracts} specifying scope, non-goals, and
  acceptance conditions. v1 preserves all v0 contracts unchanged, with
  one exception: the IVF slice's contract is \textbf{additively amended}
  through ACR-004's promotion (§9.2), expanding scope from ``input
  validation and controlled failure'' to ``input/output validation and
  controlled failure'' without modifying or removing prior scope,
  non-goals, or acceptance conditions.
\end{itemize}

\textbf{Additive OWL-level structural rigour at v1 cycle-close.} Beyond
preservation of the schema elements above, v1's cycle-close state
includes OWL-level structural axioms and metadata not expressed in v0's
YAML-only serialisation: the class-level \texttt{owl:disjointWith} axiom
between \texttt{ac:ControlledVocabularyValue} and
\texttt{ac:AppSecCoreEntity} (§3.1), ontology-header \texttt{dcterms}
metadata declared at the ontology IRI (publisher, source, issued date,
status, depiction), and per-instance \texttt{dcterms} metadata declared
at the populated-instance IRIs. These refinements are structurally
additive at the OWL level and \textbf{do not alter} any preservation
claim above: typed-instance counts are preserved (Table 2), the slice
partition is preserved, cross-slice relation types are preserved, and
per-slice contracts are preserved. The refinements have zero effect on
the apparatus's SHACL conformance verdicts under both validators (§4.2)
and zero effect on the cycle's substrate-grounding and
embedding-similarity tooling {[}4{]}.

\subsubsection{6.2 What is added across v0→v1: additive instance
population governed by the ACR
protocol}\label{what-is-added-across-v0v1-additive-instance-population-governed-by-the-acr-protocol}

The corpus expansion absorbed by the cycle reported in {[}4{]} is
reflected in v1 entirely as \textbf{additive instance population} under
the typed entity schema. The v0 → v1 entity counts below are sealed by
per-identifier audit between the v0 frozen state (released as the
\texttt{ontology-v0-frozen} artefact in the cycle's bundle, §11) and the
v1 final state (released as \texttt{ontology-v1-final}); the per-entity
verification trail is preserved in the bundle's audit-trail surface.

\textbf{Table 2.} AppSec Core v0 → v1 typed-instance counts. v0 counts
are read directly from the \texttt{ontology-v0-frozen} artefact's
instance index (\texttt{appsec-core-v0-instance-index.yaml}); v1 counts
are read directly from the same file at the \texttt{ontology-v1-final}
artefact. The v0 published surface in {[}1{]} enumerates the same 234
typed instances; v0 cycle-entry equals v0 published (no within-slice
refinement pre-iteration).

\begin{longtable}[]{@{}lrrr@{}}
\toprule\noalign{}
Entity type & v0 {[}1{]} & v1 cycle-close & Net delta \\
\midrule\noalign{}
\endhead
\bottomrule\noalign{}
\endlastfoot
ControlObjective & 70 & 75 & +5 \\
Practice & 63 & 69 & +6 \\
Mechanism & 48 & 58 & +10 \\
Artifact & 53 & 57 & +4 \\
\textbf{Total typed instances} & \textbf{234} & \textbf{259} &
\textbf{+25} \\
\end{longtable}

Slice count is preserved at 10 across v0 and v1.

\subsubsection{6.3 Attribution of the delta to ACR-protocol decisions
and within-slice
refinement}\label{attribution-of-the-delta-to-acr-protocol-decisions-and-within-slice-refinement}

The +25 delta decomposes by attribution to the cycle's worked decisions
under the ACR protocol of §7 and to a single within-slice 4-Mechanism
refinement patch. The per-source attribution below is sealed by
per-identifier audit between the \texttt{ontology-v0-frozen} and
\texttt{ontology-v1-final} bundle artefacts; the per-identifier
introduction record is preserved in the bundle's audit-trail surface
(§11).

The per-identifier introduction record (which cycle artefact-state
introduced each new entity identifier and its associated relation
declarations) is preserved in the bundle's audit-trail surface; this
section enumerates the identifiers per source and Table 3 summarises the
per-source attribution net counts.

The identifiers introduced under each attribution source are:

\begin{itemize}
\tightlist
\item
  \textbf{ACR-001} (RPR slice): \texttt{ACO-RPR-008},
  \texttt{ACO-RPR-009}, \texttt{ACO-RPR-010}; \texttt{ACP-RPR-008},
  \texttt{ACP-RPR-009}, \texttt{ACP-RPR-010}; \texttt{ACM-RPR-008},
  \texttt{ACM-RPR-009}, \texttt{ACM-RPR-010}; plus four supporting RPR
  Artifacts.
\item
  \textbf{ACR-002} (TMR slice): \texttt{ACO-TMR-008};
  \texttt{ACP-TMR-008}, \texttt{ACP-TMR-009}; \texttt{ACM-TMR-007},
  \texttt{ACM-TMR-008}.
\item
  \textbf{ACR-004} (IVF slice): \texttt{ACO-IVF-008};
  \texttt{ACP-IVF-007}; \texttt{ACM-IVF-005}.
\item
  \textbf{Within-slice refinements} (4-Mechanism patch):
  \texttt{ACM-RPR-005} (release-promotion); \texttt{ACM-SCBI-006}
  (software-composition-and-build-integrity); \texttt{ACM-SLG-005},
  \texttt{ACM-SLG-006} (secure-logging).
\end{itemize}

\textbf{Table 3.} v0 → v1 entity-count delta --- per-source attribution.

\begin{longtable}[]{@{}
  >{\raggedright\arraybackslash}p{(\linewidth - 10\tabcolsep) * \real{0.1304}}
  >{\raggedleft\arraybackslash}p{(\linewidth - 10\tabcolsep) * \real{0.1739}}
  >{\raggedleft\arraybackslash}p{(\linewidth - 10\tabcolsep) * \real{0.1739}}
  >{\raggedleft\arraybackslash}p{(\linewidth - 10\tabcolsep) * \real{0.1739}}
  >{\raggedleft\arraybackslash}p{(\linewidth - 10\tabcolsep) * \real{0.1739}}
  >{\raggedleft\arraybackslash}p{(\linewidth - 10\tabcolsep) * \real{0.1739}}@{}}
\toprule\noalign{}
\begin{minipage}[b]{\linewidth}\raggedright
Source of addition
\end{minipage} & \begin{minipage}[b]{\linewidth}\raggedleft
+CO
\end{minipage} & \begin{minipage}[b]{\linewidth}\raggedleft
+P
\end{minipage} & \begin{minipage}[b]{\linewidth}\raggedleft
+M
\end{minipage} & \begin{minipage}[b]{\linewidth}\raggedleft
+A
\end{minipage} & \begin{minipage}[b]{\linewidth}\raggedleft
Total
\end{minipage} \\
\midrule\noalign{}
\endhead
\bottomrule\noalign{}
\endlastfoot
\textbf{ACR-001} \emph{Secure Configuration Baseline Integrity} (RPR
slice ASC-09; §8.1) & +3 & +3 & +3 & +4 & 13 \\
\textbf{ACR-002} \emph{Security Requirements Lifecycle Management} (TMR
slice ASC-05; §8.2) & +1 & +2 & +2 & 0 & 5 \\
\textbf{ACR-004} \emph{Output Rendering Safety / Context-Aware Encoding}
(IVF slice ASC-07; §9.2) & +1 & +1 & +1 & 0 & 3 \\
\textbf{Within-slice refinements} (4-Mechanism patch) & 0 & 0 & +4 & 0 &
4 \\
\textbf{Total v0 → v1 net} & \textbf{+5} & \textbf{+6} & \textbf{+10} &
\textbf{+4} & \textbf{+25} \\
\end{longtable}

The decomposition reconciles bit-identically against the totals in Table
2 (+5 CO + 6 P + 10 M + 4 A = +25; v0 234 → v1 259). No entity removals
occurred between v0 and v1; identifier stability is preserved (every
entity identifier in v0's instance index is preserved in v1). The
within-slice 4-Mechanism patch is a single coherent change to existing
slice contracts: ACM-RPR-005 refines an existing ControlObjective in the
release-promotion slice without expanding scope; ACM-SCBI-006 refines
existing scope in supply-chain-and-build-integrity; ACM-SLG-005 and
ACM-SLG-006 refine the security-event-logging slice. None of the four
required ACR review under §7's threshold because none was
cross-source-convergent at the protocol's bar (≥5 independent sources
from ≥3 organisational authorities) and none was scope-amending at the
slice level --- each refines within existing contract scope and
population. Their authority for inclusion is the slice's existing
per-slice contract; the cycle's per-source curation {[}4{]} surfaced
them as instance-level refinements consistent with those contracts.

\textbf{Within-slice refinement governance discipline.} Within-slice
refinements that do not meet the ACR threshold (§7.2) are admitted under
a lighter governance discipline than ACR-level changes, with three
constraints: (a) each refinement is recorded per identifier in the
artefact's release notes and is verifiable against the bundle's
audit-trail surface (§11); (b) each refinement must remain consistent
with the slice's existing per-slice contract --- within-slice
refinements cannot amend slice scope, non-goals, or acceptance
conditions (any candidate that requires such amendment is escalated to
ACR review under §7); and (c) cumulative within-slice refinements that
exceed a defined volume threshold between major releases are escalated
to ACR-level review to guard against silent slice-scope drift. The
volume threshold (registered as future programme work: provisionally ≥5
within-slice refinements per slice or ≥20 corpus-wide between major
releases) operationalises the guard against backdoor scope amendment via
aggregation. The four within-slice Mechanism additions documented above
remain well below this provisional threshold and are recorded for
per-identifier audit at the cycle-close release.

\subsubsection{6.4 Verification trail}\label{verification-trail}

The per-source attribution in Table 3 is verifiable per identifier
against the bundle's audit-trail surface (§11), which records --- for
each of the 25 net additions --- the cycle artefact-state at which the
identifier was introduced and the corresponding relation declarations
introduced alongside it. A reader auditing the v0 → v1 transition
against the published surface of {[}1{]} can reconstruct Table 3's
per-source attribution from the audit-trail surface alone; the audit
trail is reproducible against the bundle's \texttt{ontology-v0-frozen}
and \texttt{ontology-v1-final} artefact states (§11).

\subsubsection{6.5 The empirical claim}\label{the-empirical-claim}

The empirical claim of Contribution 2 is that, at the cycle-close corpus
of thirty-one sources (a 6.2-fold source-count expansion from the
five-source first wave {[}2{]}), the v0 schema (slice partition +
populated entity types + EvidencePattern declarative class + cross-slice
relation model + per-slice contracts) admitted the corpus pressure as
additive instance population, with the entity-count delta governed by
the ACR protocol's worked decisions and a documented set of within-slice
refinements. The claim is bounded to the cycle-close corpus; whether the
schema continues to absorb additional expansion additively under future
source pressure is an empirical hypothesis carried by the cycle's
reusable workflow {[}4{]}, not a property guaranteed by the artefact
published here.

\textbf{Source-independence caveat.} The cycle-close corpus should not
be interpreted as thirty-one fully orthogonal semantic systems; several
sources belong to partially overlapping governance ecosystems. The OWASP
family --- ASVS, SAMM, DSOMM, Proactive Controls, and the OWASP Top-10
specialisations --- shares editorial framing and convergent vocabulary
even where individual documents target distinct surfaces. The NIST
family (SSDF, SP 800-53, AI 100-2, AI RMF) is similarly co-authored
within a single organisational authority. The ACR protocol's
\texttt{≥3\ organisational\ authorities} qualifier (§7.2 Condition 1)
mitigates but does not eliminate this dependency: it ensures convergence
is not carried by a single authority's editorial framing, but partial
vocabulary co-evolution within and across authorities remains a property
of the corpus. The substantive claim of C2 is therefore that the v0
schema absorbed the corpus's \emph{partially-overlapping} multi-source
pressure as additive instance population, not that it absorbed
thirty-one independent semantic systems.

\textbf{Regulatory horizontal sources in the cycle-close corpus.} The
corpus includes horizontal regulatory frameworks alongside the
AppSec-canonical and AI/ML extensions: PCI DSS v4.0.1 and PCI SSLC v1.1
(PCI SSC); EU CRA, EU NIS2, EU GDPR, EU DORA (EU regulatory wave); HIPAA
Security Rule (US HHS). Per-source claim density to v1 typed entities
varies substantially across these sources --- PCI DSS contributes
high-density mappings (e.g., 104 claims to the RPR slice's
ACR-001-promoted ControlObjective set alone; substantial mappings across
IAT, IVF, SCBI), consistent with PCI DSS's heavy engineering-substance
loading. EU horizontal regulations (CRA, NIS2, GDPR, DORA) and HIPAA
contribute lower-density mappings, consistent with their
governance-substance loading: bulk of these regulations' content is
governance substance (board-level reporting, vendor risk management,
regulatory examination response, organisation-level policy authority,
workforce training programmes) routed by design to the Manual's prose
surface (§1.3); the engineering substance fraction that \emph{does}
normalise to v1 typed entities (technical safeguards, vulnerability
management dispositions, incident detection mechanisms, supply-chain
integrity controls) maps coherently into the existing slice partition
without requiring schema disruption. This empirical pattern across the
regulatory subset is consistent with the engineering-substance /
governance-substance design split declared in §1.3. The
framework-horizontal stress test --- explicit decomposition against each
regulatory framework with per-framework engineering-substance /
governance-substance fractioning --- is registered as future programme
work; the present paper reports the cycle-close corpus's absorption
pattern at the aggregate level.

A potential failure mode of the schema-preservation claim --- what would
have counted as failure under the cycle's measurement --- is enumerated
explicitly: a slice added or removed; a populated entity type added; a
cross-slice relation added or modified; a per-slice contract
restructured (rather than additively amended); an ACR candidate that the
protocol's threshold could not bound under any of (i)
within-existing-slice-amendment, (ii)
within-existing-population-refinement, or (iii) explicit non-admission
with identifier consumption. None of these failure modes occurred at the
cycle-close corpus. The schema-preservation claim is anchored to the
absence of these failure modes, not to a generic claim of stability.

\textbf{Anti-tautological defense.} The ACR protocol's Condition 4
(§7.2) excludes partition-level restructuring --- a sibling-slice
repartition or addition of a new populated entity type --- from the
protocol's scope; such a candidate would require a separate
partition-level governance exercise. This design choice makes ``no slice
repartition'' a structural property of the protocol rather than an
empirical observation: the protocol cannot, by construction,
\emph{produce} partition-level disruption from within. The
schema-preservation evidence of Contribution 2 is therefore:
\emph{across the cycle-close thirty-one-source corpus and the four
worked ACR decisions, no candidate's substance required partition-level
restructuring to be normalized}. If a candidate had emerged whose
substance could not be admitted under any of the in-scope ACR
dispositions (within-existing-slice-amendment,
within-existing-population-refinement, or explicit non-admission with
identifier consumption), the protocol would have flagged it as
out-of-scope, requiring a separate partition-level governance exercise
--- a path the cycle did not need to invoke. The substantive claim is
therefore the \emph{absence of partition-level pressure} under the
cycle-close corpus, not the protocol's inability to handle such pressure
if it arose. Whether future expansion produces a candidate that requires
partition-level governance is an empirical question carried by the
cycle's reusable workflow, not foreclosed by the protocol's design.

\textbf{Out-of-scope candidate register (the audit-trail counterpart).}
To make ``no partition-level pressure emerged'' verifiable rather than a
structural absence, the protocol's audit-trail surface (§11) records ---
for every candidate the cycle considered --- its disposition under §7.2
Condition 4: in-scope (admitted to one of the three in-scope
dispositions: promotion, within-population refinement, or non-admission
with identifier consumption) or out-of-scope (classified as requiring
partition-level governance, with the reason recorded). At cycle close,
the out-of-scope register contains \textbf{N = 0 entries}: no candidate
the cycle considered was classified as out-of-scope of the protocol's
Condition 4. This zero-entry register is the verifiable counterpart of
the substantive claim --- the absence is empirical, not constructive.
The register is preserved in the bundle's audit-trail surface alongside
the four in-scope worked decisions of §§8--9.

\textbf{Relation to conservative extension in the ontology-engineering
literature.} The schema-preservation claim of Contribution 2 is
\emph{weaker} than conservative extension stricto sensu {[}30, 31{]}: it
does not claim preservation of entailments over the v0 vocabulary, but
claims that (i) the v0 TBox (slice partition + populated entity type
declarations + cross-slice relation type set) is preserved unchanged in
v1; (ii) the cross-slice relation type system is unchanged; and (iii)
the v1 ABox contains the v0 ABox under identifier stability (every v0
entity identifier is preserved in v1 with its full property set carried
forward). The relation to module-theoretic ontology evolution
(locality-based module extraction; formal conservative-extension
verification under the Σ-conservative-extension definition) is
registered as future programme work; the claim reported here is the
weaker structural-and-identifier preservation under additive ABox
population.

\begin{center}\rule{0.5\linewidth}{0.5pt}\end{center}

\subsection{7. The AppSec Core Change Request (ACR)
protocol}\label{the-appsec-core-change-request-acr-protocol}

\subsubsection{7.1 What an ACR is}\label{what-an-acr-is}

An \textbf{AppSec Core Change Request (ACR)} is a candidate change to
the bounded ontology submitted under the protocol described in this
section. Candidates may be:

\begin{itemize}
\tightlist
\item
  a candidate new ControlObjective, Practice, Mechanism, or Artifact;
\item
  a candidate new cross-slice relation;
\item
  a candidate scope amendment to a slice's per-slice contract (additive
  amendment within bottom-up derivation discipline);
\item
  a candidate non-admission decision recording that a candidate's
  substance does not meet the protocol's threshold and is held outside
  the ontology, with the reason recorded.
\end{itemize}

An ACR receives a unique identifier of the form \texttt{ACR-NNN},
assigned at submission time, and consumes that identifier permanently
regardless of outcome (promotion, non-admission, or any future
re-evaluation under a different threshold). The identifier-consumption
discipline preserves a complete record of every governance decision the
protocol has produced, against which future submissions can be audited.

\subsubsection{7.2 The four-condition promotion
threshold}\label{the-four-condition-promotion-threshold}

The protocol's threshold for promotion is stated as four conditions, all
of which must be met for a candidate to be admitted to the ontology:

\textbf{Condition 1 --- Multi-source convergence.} The candidate's
substance appears across \textbf{at least five independent sources drawn
from at least three distinct organisational authorities}. The
five-source threshold is calibrated against the cycle's first-wave
bounding test {[}2{]}: at five sources from independent authorities,
convergence on the same substance establishes that the candidate is not
idiosyncratic to a single team's interpretation. Lower thresholds (three
or four sources) are detection-level signals that may motivate further
investigation but do not on their own meet the promotion bar. The
``three organisational authorities'' qualifier prevents a candidate from
being inflated by multiple publications of a single authority counting
as multiple sources. The five-source / three-authorities threshold is
\textbf{governance-calibrated rather than mathematically optimal}: it
operationalises the trade-off between admitting clusters carried by
single-pattern editorial framing across few authorities (a lower
threshold would do this) and excluding substantively important but
distinctively-framed concepts (a higher threshold would do this), at a
heuristic balance point anchored to the first-wave bounding test.
Sensitivity analysis under alternative thresholds (three-source minimum;
seven-source higher bar) is registered as future work (§10).

\textbf{Condition 2 --- Multi-method convergence.} The candidate is
detected by independent normalization signals during the cycle's
re-application of the refined method {[}4{]}: a domain-first signal
reading the source's own structural hierarchy to place the candidate at
a coherent slice, a sentence-embedding-similarity signal returning a
coherent top-anchor cluster across the convergent sources, and an
explicit override-mechanism record where the author's curated decision
concurs with or overrides the method-converged signal under
audit-traceable discipline. All three signals are required;
method-single-signal detection does not on its own meet Condition 2.

\textbf{Condition 3 --- Practitioner-Manual content backing.} The
candidate's substance has a counterpart in the prescriptive practitioner
Manual that complements the bounded ontology in the programme
architecture {[}5{]}. Where the Manual does not yet cover the substance,
the cycle's gap-detection {[}4{]} prompts the Manual's author to fill
the missing content before the candidate proceeds; ontology promotion
never precedes Manual-prose authorship of the substance. The Manual's
authoring is \textbf{iteratively responsive to ACR pressure rather than
antecedent to it}: a candidate surfacing from source convergence may
trigger gap-detection against the Manual, which the Manual's authoring
process then resolves; ontology promotion under §7 follows Manual-prose
authorship, but the Manual itself evolves under cycle pressure (the
bidirectional Manual-validation discipline is reported in {[}5{]}). The
condition's purpose is to keep ontology change tied to authored
prescriptive content rather than to source-side convergence alone.

\textbf{Condition 4 --- Slice fit.} The candidate fits within an
existing slice's per-slice contract (its scope, non-goals, acceptance
conditions), or motivates a slice-contract amendment whose acceptance
under the bottom-up derivation discipline can be defended without
restructuring sibling slices. Restructuring a sibling slice or
repartitioning the ten-slice partition is outside the protocol's scope;
such a candidate would be a partition-level decision requiring its own
(separate) protocol exercise.

\subsubsection{7.3 Symmetric application: promotion or
non-admission}\label{symmetric-application-promotion-or-non-admission}

The protocol is symmetric: under the same four-condition threshold, a
candidate may be \textbf{promoted} (admitted to the ontology, with the
promoted entities entered at the populated graph) or \textbf{not
admitted} (held outside the ontology, with the reason recorded against
whichever condition failed and the identifier consumed). Symmetric
application of the same conditions to both outcomes is the protocol's
mechanism for preventing retrofit: if non-promotion never occurred under
the protocol, the threshold would be unfalsifiable. Section 9.1 reports
the cycle's first non-admission decision (ACR-003), which alongside the
three promotions of §§8.1, 8.2, and 9.2 provides initial evidence that
the threshold produces falsifiable governance decisions; near-miss
non-admission cases that would test single-condition marginal failures
are registered as future work (§10.6).

\subsubsection{7.4 The protocol's evidence
trace}\label{the-protocols-evidence-trace}

For each ACR (admitted or not), the protocol records: (a) the
candidate's substance and proposed entities or scope amendment; (b) the
per-condition evidence (the convergent sources for Condition 1; the
per-stage method signals for Condition 2; the Manual content reference
for Condition 3; the slice contract reference and amendment proposal if
any for Condition 4); (c) the disposition (promotion, non-admission);
(d) where promoted, the assigned entity identifiers and the
slice-contract delta if any; (e) where not admitted, the reason against
the failed condition. The evidence trace is recorded against the
artefact's release notes at the cycle-close release tag, with each ACR's
per-condition evidence verifiable against the cycle's published
per-source records {[}4{]}.

\begin{center}\rule{0.5\linewidth}{0.5pt}\end{center}

\subsection{8. Worked decisions on candidates carried forward from the
v0
baseline}\label{worked-decisions-on-candidates-carried-forward-from-the-v0-baseline}

\subsubsection{8.1 ACR-001 --- Secure Configuration Baseline Integrity
(promoted)}\label{acr-001-secure-configuration-baseline-integrity-promoted}

\textbf{Substance.} The integrity of secure configuration baselines as a
first-class engineering practice in the
release-promotion-controlled-rollout-and-rollback-readiness slice
(ASC-09 / RPR family): the existence of an authoritative secure-baseline
configuration for security-relevant runtime parameters, the integrity of
security-relevant configuration values against unauthorized override,
and the discipline of baseline review with exception visibility and
change discipline. Configuration baseline integrity governs what is
deployed to runtime and is therefore canonically placed in the
release-promotion slice, which carries the release-time controls that
gate baseline integrity at promotion and rollback boundaries.

\textbf{v0-era status.} The substance was already present in the v0
Manual mapping {[}1, §10{]} as \texttt{CFG-001→007} (``secure
configuration baselines\ldots{} documented in the corpus in 2023'') with
anteriority noted relative to ASVS v5.0's 2025 formalisation of
\texttt{secure\_configuration\_baseline}. v0's published surface
reflected the substance under the Manual-mapping label without
first-class ControlObjective status under an explicit governance
protocol.

\textbf{Per-condition evidence.} - \emph{Condition 1 (multi-source
convergence at ≥5 sources / ≥3 organisational authorities).} At
substrate cycle close, \textbf{27 of the 31 corpus sources} carry at
least one normalized claim targeting one of the three promoted
ControlObjectives (ACO-RPR-008/009/010); the 27 sources span \textbf{8
distinct organisational authorities} --- NIST, OWASP, MITRE, CIS, PCI
SSC, SAFECode, EU regulatory wave, and US HHS --- substantially
exceeding the §7 four-condition threshold of ≥5 sources from ≥3
authorities. Configuration baseline integrity is a near-universal AppSec
topic across the corpus; the per-source claim distribution concentrates
in NIST SP 800-53, MITRE CAPEC, PCI DSS, OWASP SAMM, CIS Controls, OWASP
ASVS, OWASP DSOMM, and NIST SSDF (top eight by claim count), with the
remaining 19 sources carrying smaller per-source claim sets. The full
per-source claim distribution is preserved in the substrate's per-pilot
configuration files at the cycle-close release tag (see {[}4 §15.2{]}).
- \emph{Condition 2 (multi-method convergence).} The cycle's
re-application of the refined method {[}4 §4{]} returns convergent
signals across the structural-hierarchy stage (the source-side ancestors
place the substance within the RPR slice), the embedding-similarity
stage (top-anchor cluster coheres across the convergent sources), and
the override-mechanism record (no override required against the
method-converged signal). - \emph{Condition 3 (practitioner-Manual
content backing).} The substance is present in the practitioner Manual
under the release-promotion and configuration-management chapters of the
Manual's prose layer; the Manual's coverage of secure configuration
baseline substance has been authored at the Manual-validation cycle's
documented state {[}5{]}. - \emph{Condition 4 (slice fit).} The
candidate fits within the existing RPR slice's per-slice contract; no
slice-contract amendment is required.

\textbf{Disposition.} Promoted. The substance is admitted to the
ontology as three new ControlObjectives entered at the populated graph
for v1: \textbf{ACO-RPR-008} (Secure Configuration Baseline Integrity),
\textbf{ACO-RPR-009} (Security-Relevant Configuration Integrity And
Override Control), and \textbf{ACO-RPR-010} (Baseline Review, Exception
Visibility And Change Discipline), together with their associated
Practice and Mechanism chains.

\subsubsection{8.2 ACR-002 --- Security Requirements Lifecycle
Management
(promoted)}\label{acr-002-security-requirements-lifecycle-management-promoted}

\textbf{Substance.} The discipline of managing security requirements as
living artefacts across the application's development cycle:
requirements creation, requirements review, requirements traceability
across sprint boundaries, requirements deprecation, and
requirements-versus-threat coherence under change. Distinct from the
threat-side substance the TMR slice (ASC-05) already carries (threat
identification, threat-and-risk traceability to mitigation, mitigation
prioritisation), the candidate anchors the
\textbf{requirements-as-living-artefacts} dimension of the TMR slice.

\textbf{v0-era status.} The substance was identified at v0 closure as a
deferred candidate; v0's published protocol thresholds were not applied
symmetrically across the v0-era candidate set {[}1{]}, and ACR-002 sat
at v0 closure in a ``candidate, not yet decided under protocol'' state.
Security-requirements lifecycle management is so foundational to AppSec
that its absence as a first-class ControlObjective set in v0 deserves
explanation. The \emph{substance} was already present in the
practitioner Manual's prose layer at v0 publication (Manual mapping
label \texttt{REQ-LCM-*} per {[}1, §10{]}); what ACR-002 promotes to v1
is \textbf{first-class typed-entity status under an explicit governance
protocol}, not the substance itself. The promotion is therefore an
\emph{artefact-formalization step} --- moving substance from Manual
prose into the bounded ontology under the ACR protocol's four-condition
discipline --- rather than a discovery of new substance unaccounted for
in the v0-era programme.

\textbf{Per-condition evidence.} - \emph{Condition 1.} At substrate
cycle close, \textbf{21 of the 31 corpus sources} carry at least one
normalized claim targeting one of the five promoted entities
(ACO-TMR-008 + ACP-TMR-008/009 + ACM-TMR-007/008); the 21 sources span
\textbf{9 distinct organisational authorities} --- NIST, OWASP, MITRE,
CIS, PCI SSC, SAFECode, SLSA, EU regulatory wave, and US HHS ---
exceeding the §7 four-condition threshold of ≥5 sources from ≥3
authorities. Top backing sources by claim count include NIST SP 800-53,
OWASP SAMM Security Requirements practice {[}16{]}, PCI DSS, NIST SSDF
{[}9{]} practice group PS, CIS Critical Security Controls, and OWASP
ASVS {[}10{]}. The full per-source claim distribution is preserved in
the substrate's per-pilot configuration files at the cycle-close release
tag (see {[}4 §15.2{]}). - \emph{Condition 2.} Multi-stage method
convergence: the structural-hierarchy stage places the candidate within
the TMR slice; the embedding-similarity stage returns coherent
top-anchor cluster across the convergent sources; the override-mechanism
record concurs with the method signal. - \emph{Condition 3.} Substance
present in the practitioner Manual under the threat-modeling chapter of
the Manual's prose layer. - \emph{Condition 4.} The candidate motivates
an additive amendment to the TMR slice's per-slice contract scope,
declaring the requirements-as-living-artefacts dimension as within scope
alongside the existing threat-side substance. The amendment is additive:
prior scope, non-goals, and acceptance conditions are preserved.

\textbf{Substantive question raised by the candidate.} \emph{Does
ACR-002 duplicate scope already in the TMR slice?} The TMR slice at v0
closure carries the threat-side substance --- threat identification,
threat-and-risk traceability, mitigation prioritisation --- with
security requirements treated as the \emph{output} of the
threat-modeling activity rather than as a governance discipline in their
own right. The candidate's substance is the
\emph{requirements-as-lifecycle-artefact} dimension, distinct from the
threat-side dimension v0 already captured. Three considerations support
the additive amendment rather than collapse into the existing
threat-side anchor: the convergent sources surface the candidate
substantively as requirements lifecycle (a temporal-and-traceability
dimension distinct from the threat-side dimension); v0's slice-contract
before the amendment did not declare a primary anchor for the
requirements-as-living-artefacts dimension; and the cycle's per-stage
method signals support the additive amendment's slice fit without
restructuring the threat-side substance. The promoted ControlObjective
therefore extends the slice along an orthogonal dimension to the
existing threat-side substance, rather than duplicating it.

\textbf{Disposition.} Promoted. The substance is admitted to the
ontology as one new ControlObjective + two new Practices + two new
Mechanisms in the TMR slice, together with the additive amendment to the
slice's per-slice contract scope.

\begin{center}\rule{0.5\linewidth}{0.5pt}\end{center}

\subsection{9. Worked decisions on candidates surfaced by cycle
re-application}\label{worked-decisions-on-candidates-surfaced-by-cycle-re-application}

\subsubsection{9.1 ACR-003 --- Multi-anchor adjacency candidate batch
(not
admitted)}\label{acr-003-multi-anchor-adjacency-candidate-batch-not-admitted}

\textbf{Substance.} A batch of approximately thirty-five candidate items
distributed across four sub-clusters surfaced when the cycle's
re-application of the refined method {[}4{]} returned elevated
embedding-similarity scores against multiple existing ControlObjectives
without converging on a single primary anchor. The batch was registered
as a single ACR candidate to evaluate whether the cycle's
similarity-stage signal had detected genuinely new ontological substance
or had surfaced items already covered at appropriate adjacency by
existing entities.

\textbf{Per-condition evidence.} - \emph{Condition 1.} Source coverage
across the cycle's substrate met the multi-source threshold at the batch
level; however, the batch's items as a coherent cluster did not converge
on a single substantive concept at the candidate level --- items
distributed across multiple existing ControlObjectives without a
unifying concept that the protocol could promote. - \emph{Condition 2.}
Multi-stage method signals were \emph{partially} aligned: the
embedding-similarity stage flagged elevated scores; the
structural-hierarchy stage did not return a coherent slice signal that
would have anchored the batch to a single substantive promotion. The
override-mechanism record carries the author's evaluation that the
batch's items, on closer review, were either already covered by existing
entities at appropriate adjacency or were items whose substantive
content fell outside the bounded ontology's scope and into the
practitioner Manual's prose layer. - \emph{Condition 3.} Where the
batch's items had practitioner-Manual content backing, the backing was
already covered by existing ontology entities at appropriate adjacency,
with no residual content motivating new substance. - \emph{Condition 4.}
No coherent slice fit emerged --- the batch's items, as a candidate
cluster, did not motivate either an existing-slice fit or a defensible
slice-contract amendment.

\textbf{Disposition.} Not admitted. The protocol's threshold is not met:
items are either already covered at appropriate adjacency by existing
entities, or are in scope for the practitioner Manual rather than for
the bounded ontology. The non-admission is reported as an active result
of the protocol --- the protocol's symmetric application produces
non-promotion under the same conditions as it produces promotion, and
the cycle's record of the non-admission is itself evidence that the
threshold is testable rather than retrofitted.

The ACR-003 identifier is \textbf{consumed}: future cycles will not
re-use the identifier for a different candidate, even if substance later
resurfaces under different framing.

\subsubsection{9.2 ACR-004 --- Output Rendering Safety / Context-Aware
Encoding
(promoted)}\label{acr-004-output-rendering-safety-context-aware-encoding-promoted}

\textbf{Substance.} Output rendering safety and context-aware encoding
for content destined for downstream interpretation (HTML rendering, SQL
execution, command-interpreter execution, log-injection contexts, and
analogous surfaces). Distinct from input validation alone --- the input
may be valid in form yet still produce unsafe output if
context-appropriate encoding is omitted at the rendering boundary. The
substance addresses the output side of the input-and-output-validation
slice's substrate.

\textbf{v0-era status.} As with ACR-002 (§8.2), output rendering safety
/ context-aware encoding is so foundational to AppSec that its absence
as first-class typed substance in v0 deserves explanation. The substance
was present in the practitioner Manual's validation chapter at v0
publication, and v0's input-validation slice carried adjacent
ControlObjectives covering input boundary discipline and controlled
failure under hostile input --- but the \emph{output side} of the
validation slice (context-aware encoding at the rendering boundary) was
not first-class typed substance under explicit governance. ACR-004 is
therefore a \emph{first-class promotion of pre-existing Manual
substance} under the ACR protocol's discipline, with an additive scope
amendment expanding the slice's contract from ``input validation and
controlled failure'' to ``input/output validation and controlled
failure''; it is not a discovery of new substance unaccounted for in the
v0-era programme.

\textbf{Per-condition evidence.} - \emph{Condition 1.} At substrate
cycle close, \textbf{8 of the 31 corpus sources} carry at least one
normalized claim targeting one of the three promoted entities
(ACO-IVF-008 + ACP-IVF-007 + ACM-IVF-005); the 8 sources span \textbf{3
distinct organisational authorities} --- MITRE (via CAPEC {[}13{]}, CWE
{[}14{]}, and ATLAS), OWASP (via ASVS V5 {[}10{]}, DSOMM {[}17{]},
Proactive Controls, and OWASP Top-10 for LLMs), and NIST (via SP
800-53). The candidate hits the §7 four-condition threshold at exact
minimum (3 organisational authorities); explicit disclosure of the
minimum-threshold case is reviewer-defensible --- the cycle ran close to
but met the threshold, demonstrating the protocol's testability rather
than retrofitting. The per-source claim distribution is preserved in the
substrate's per-pilot configuration files at the cycle-close release tag
(see {[}4 §15.2{]}). - \emph{Condition 2.} Multi-stage method
convergence: the structural-hierarchy stage places the candidate within
the IVF slice with the slice-contract amendment proposed; the
embedding-similarity stage returns coherent top-anchor cluster across
the convergent sources; the override-mechanism record concurs with the
method signal. - \emph{Condition 3.} Substance present in the
practitioner Manual under the validation chapter of the Manual's prose
layer. - \emph{Condition 4.} The candidate motivates an additive
amendment to the IVF slice's per-slice contract scope, expanding from
``input validation and controlled failure'' to ``input/output validation
and controlled failure''. The amendment is additive: prior scope (input
validation, controlled failure under hostile or malformed input) is
preserved; the candidate adds the output-side substance.

\textbf{Disposition.} Promoted. The substance is admitted to the
ontology as one new ControlObjective + one new Practice + one new
Mechanism in the (now-expanded) IVF slice, together with the additive
amendment to the slice's per-slice contract scope.

\subsubsection{9.3 The four-decision symmetric
application}\label{the-four-decision-symmetric-application}

The four-decision span --- ACR-001 promoted (§8.1), ACR-002 promoted
(§8.2), ACR-003 not admitted (§9.1), ACR-004 promoted (§9.2) --- is the
evidence supporting Contribution 3. Three promotions and one
non-admission decided under the same four-condition threshold provide
initial evidence that the protocol produces falsifiable governance
decisions: the same conditions that produced promotion under ACR-001,
ACR-002, and ACR-004 produced non-admission under ACR-003. The
protocol's threshold is therefore testable rather than retrofitted; the
symmetric application is the record against which future ACR decisions
can be audited for consistency with the published threshold, and §10.6
registers the empirically grounded near-miss non-admission case as
future work to close the discriminative-power evidence against
single-condition marginal failures.

\begin{center}\rule{0.5\linewidth}{0.5pt}\end{center}

\subsection{10. Limitations}\label{limitations}

\subsubsection{10.1 Independent-research
scope}\label{independent-research-scope}

This artefact and its governance protocol are published as independent
research {[}4 §8.3{]}. The following limits are declared honestly and
registered as future work:

\begin{itemize}
\tightlist
\item
  \emph{No subject-matter-expert panel review of individual mappings.}
  Per-instance mapping decisions and the per-condition evidence for each
  of the four ACR worked decisions are author-curated under the cycle's
  documented discipline {[}4{]}. Subject-matter-expert review of
  individual mappings is not part of the artefact's scope; future
  iterations of the cycle may incorporate such review and revise the
  artefact accordingly.
\item
  \emph{No inter-rater reliability measurement.} The per-instance
  population of v1, the within-slice refinement attributions of §6.3,
  and the four ACR worked decisions are single-evaluator
  categorisations. Inter-rater reliability validation of these
  categorisations is not part of this paper's scope.
\item
  \emph{No peer-validated submission of individual mappings.} The
  cycle's per-source-into-Reference mapping skin (the OLIR-compatible
  publication form reported in {[}4{]}) is form-level only; peer
  submission to a programme such as NIST OLIR or to a related
  cross-mapping community is not part of this paper's scope and is
  registered as future work.
\end{itemize}

\subsubsection{10.2 Schema-preservation claim is bounded to the
cycle-close
corpus}\label{schema-preservation-claim-is-bounded-to-the-cycle-close-corpus}

Contribution 2's schema-preservation claim is bounded to the
thirty-one-source corpus at the cycle's close. Future source expansion
may surface ACR pressure that the protocol's current threshold does not
anticipate, or may surface substance that the schema cannot absorb
additively under existing slice contracts. The cycle's reusable workflow
{[}4{]} is the response: re-invocation of the cycle's method on a new
source is the discipline by which future expansion is integrated, with
the protocol of §7 governing whether new ontology change is warranted.
The artefact published here does not claim future-source bounding.

\subsubsection{10.3 The bidirectional Manual-validation cycle is
downstream}\label{the-bidirectional-manual-validation-cycle-is-downstream}

The bidirectional validation cycle that exercises v1 against the
practitioner Manual's prose layer --- the Manual covering v1's instance
population in prose, the Manual's coverage of substance routed outside
v1's bounds, the edit-versus-traceback decisioning for each of the four
ACR worked decisions --- is reported in {[}5{]}. Misfit surfaced by
{[}5{]} may iterate the artefact reported here. The artefact's release
at the cycle-close release tag is the artefact's state at the design
science cycle's close; revision under {[}5{]}'s validation cycle would
produce a new release tag.

\subsubsection{10.4 EvidencePattern population is
deferred}\label{evidencepattern-population-is-deferred}

EvidencePattern is preserved as a declarative class in v1 (§2.2) without
additional populated instances at this release. The v0 supporting index
of evidence patterns {[}1, §6{]} is preserved; population of the
EvidencePattern class as first-class typed instances is registered as
future programme work. Downstream consumers requiring deterministic
evidence-pattern verification can use the v0 supporting index against
the v1 populated graph at the cycle-close release tag; first-class
typed-instance population would extend the artefact's
machine-verifiability surface.

\subsubsection{10.5 Versioning policy}\label{versioning-policy}

The artefact follows a semantic-versioning policy aligned with the ACR
protocol's governance scope:

\begin{itemize}
\tightlist
\item
  A \textbf{major version increment} (v1 → v2) reflects a
  partition-level change: addition or removal of a slice, addition of a
  populated entity type, modification of the cross-slice relation type
  set, or restructuring of an existing slice's contract beyond additive
  amendment. Such changes are out of scope of the ACR protocol (§7.2
  Condition 4) and require a separate partition-level governance
  exercise.
\item
  A \textbf{minor version increment} (v1 → v1.1, v1.2) reflects
  ACR-promoted additive change under the protocol of §7: a promoted
  ControlObjective / Practice / Mechanism / Artifact set, an additive
  amendment to an existing slice's contract scope, or a new ACR
  non-admission decision recorded against the protocol's threshold.
\item
  A \textbf{patch version increment} (v1.0 → v1.0.1, v1.0.2) reflects
  within-slice refinements admitted under the lighter governance
  discipline of §6.3 (per-identifier audit + slice-contract consistency
  + cumulative-volume threshold), or non-substantive editorial revisions
  (metadata, formatting, controlled-vocabulary fixes) that do not alter
  populated content or relation declarations.
\end{itemize}

v1 as published here is the artefact's first OWL/SHACL formal release;
v0 {[}1{]} was YAML-only. The v0 → v1 transition is therefore an
artefact-formalization step that decomposes into the additive instance
population documented in §6 (minor-version equivalent: ACR-promoted
additions) plus the within-slice 4-Mechanism refinement (patch-version
equivalent). Future evolution under the policy above is the artefact's
evolution path; partition-level revisions remain subject to separate
programme-lead governance outside the ACR protocol's scope.

\subsubsection{10.6 Discriminative power against near-miss non-admission
is not empirically
demonstrated}\label{discriminative-power-against-near-miss-non-admission-is-not-empirically-demonstrated}

The protocol's discriminative power under the four-condition threshold
is established by the protocol's specification (§7.2), and the symmetric
application across the four worked decisions of §§8--9 (three
promotions, one non-admission) is reported as evidence that the protocol
is testable rather than retrofitted. The four worked decisions do
\textbf{not} include a \emph{near-miss non-admission} --- a candidate
that failed a single condition by a marginal amount (for example, a
candidate meeting Conditions 2, 3, and 4 but missing Condition 1 by
carrying convergent substance across only four sources or two distinct
organisational authorities, or meeting Conditions 1, 2, and 3 but
failing Condition 4 by motivating a slice-contract amendment that could
not be defended without sibling-slice restructuring). ACR-003's
non-admission was a \emph{framing-level failure} (batch of approximately
thirty-five candidate items without a unifying substantive concept
across the batch, §9.1), categorically different from a single-condition
marginal failure. ACR-004's promotion at exact minimum threshold (8
sources / 3 organisational authorities, §9.2) demonstrates that the
protocol accepts close to the boundary but does not demonstrate that it
rejects close to the boundary. The protocol's discriminative power
against single-condition marginal failures is therefore established by
the threshold's specification, not by empirical worked decisions; an
empirically grounded near-miss non-admission case is registered as
future programme work under the cycle's reusable workflow {[}4{]}, to
surface in subsequent cycle iterations when a candidate of that profile
arises.

\textbf{Conservative-bias-by-design.} The four-condition threshold is
calibrated as \textbf{conservative-by-design}: it favours a low
false-positive admission rate (promoting candidates that should not be
admitted) at the operational cost of a higher false-negative
non-admission rate (declining candidates that could perhaps be admitted
but fail one condition by a marginal amount). This bias is the
defensible governance posture for a bounded ontology where every
admission carries downstream cost (apparatus regeneration; MCP grounding
refresh; Manual cross-reference update; cross-paper bilateral
consistency maintenance across the programme's papers); under-admission
can be reversed by re-promotion under a future cycle, whereas
over-admission contaminates the ontology's bounded surface and the
downstream tooling that consumes it. The asymmetric discriminative-power
profile of §10.6 is therefore not an empirical gap but a design feature
of the protocol's calibration.

\subsubsection{10.7 External-tool validation (OOPS! and
FOOPS!)}\label{external-tool-validation-oops-and-foops}

AppSec Core v1 passes the OntOlogy Pitfall Scanner OOPS! {[}35{]} with
zero Critical and zero Important pitfalls. Two Minor pitfalls remain as
documented design choices. The first is OOPS! P04 (unconnected ontology
element), which affects the abstract grouping superclass
\texttt{ac:ControlledVocabularyValue}: by design this class carries no
direct instance binding (its six controlled-vocabulary subclasses are
the instance-bearing classes --- the taxonomy-root pattern), and it is
connected at class-level via \texttt{rdfs:subClassOf} from those six
subclasses plus the \texttt{owl:disjointWith\ ac:AppSecCoreEntity} axiom
declared in §3.1. The second is OOPS! P22 (different naming
conventions), which reflects a deliberate convention split:
\texttt{snake\_case} for schema-derived properties (mirroring canonical
YAML keys to keep the OWL surface trivially round-trippable to/from the
YAML schema source) and \texttt{camelCase} for OWL-generated derived
properties (following the established RDF \texttt{hasX} / \texttt{isXOf}
idiom). Unifying the conventions would require either breaking YAML key
fidelity or breaking RDF community readability; both alternatives carry
higher cost than the documented split.

AppSec Core v1 scores 13/15 on the FAIR Ontology Pitfall Scanner FOOPS!
{[}36{]} binary baseline. The two remaining gaps are gated on external
prerequisites: PURL1 (persistent URL) pending registration of a
\texttt{w3id.org} redirect for the ontology namespace, registered as
future programme work in the cycle's reusable workflow {[}4{]}; and the
OM3 binary check (detailed reusability metadata) pending DOI assignment
at bundle deposit (five of six OM3 sub-fields are already present at the
cycle-close release --- publisher, logo, status, source, issued --- and
the missing sub-field is \texttt{doi}). The FOOPS! score will lift to
14/15 when the deposit DOI is assigned and added to the ontology header,
and to 15/15 if the \texttt{w3id.org} persistent-IRI redirect is
subsequently registered.

\begin{center}\rule{0.5\linewidth}{0.5pt}\end{center}

\subsection{11. Reproducibility}\label{reproducibility}

\subsubsection{11.0 The deposited bundle}\label{the-deposited-bundle}

This paper's published artefacts are deposited as \textbf{a single
bundle} at a public archive (DOI to be assigned at deposit; interim
citation form: \emph{cycle-close bundle, archive DOI pending public
release}). Within the bundle, the artefacts referenced throughout this
paper are identified by stable artefact identifiers ---
\texttt{ontology-v1-final} (the OWL 2 DL ontology + populated graph;
§11.1), \texttt{apparatus-shacl-pyshacl-v3} (the SHACL apparatus;
§11.2), \texttt{appsec-core-embeddings-v1.1} (the cycle's
embedding-similarity tooling, used by the cycle reported in {[}4{]}) ---
together with the audit-trail surface that records per-identifier
introduction provenance for the v0 → v1 entity-count attributions of §6
(Tables 2--3).

Each artefact within the bundle is SHA-256-pinned at the bundle's
manifest. Where this paper cites ``the cycle-close release'', the
reference is to the bundle deposit; readers verify the SHA-256 pins
reported per artefact below against the manifest's hashes, without
requiring access to the cycle's working repositories.

\subsubsection{11.1 OWL 2 DL ontology}\label{owl-2-dl-ontology}

The ontology export \texttt{ontology-v1-final} is published in Turtle
(\texttt{.ttl}) format within the bundle, SHA-256-pinned at the
manifest. A reader applying a standard OWL 2 DL reasoner (e.g., HermiT,
Pellet, OpenLLET) to the populated graph reproduces the class-hierarchy
and property-domain/range checks reported in §3.

\subsubsection{11.2 SHACL apparatus}\label{shacl-apparatus}

The apparatus \texttt{apparatus-shacl-pyshacl-v3} is published as two
complementary shape files within the bundle: the schema-derived ontology
shapes file (\texttt{appsec-core-v0-shapes.ttl}) and the
consumer-conformance shapes file
(\texttt{consumer-conformance-shapes.ttl}), together with the build
script (\texttt{build\_shacl.py}) for the schema-derived shapes and its
documented build environment. Re-running the build script under the
documented environment produces a deterministic SHACL output
reproducible against the released schema-derived shapes file (isomorphic
up to blank node labeling under the RDFLib serialisation defaults of the
documented environment); SHA-256 pinning at the bundle's manifest is the
canonical reference for all three files.

\subsubsection{11.3 Validation pair}\label{validation-pair}

The validation pair is \texttt{pyshacl} version 0.31.0 with
\texttt{rdflib} 7.6.0 (the SHACL Core reference implementation used here
as canonical validation oracle, community-maintained against the W3C
SHACL Recommendation rather than W3C-endorsed) plus the in-house
bounded-subset validator at parity for the apparatus's six ontology
shapes. A reader applying \texttt{pyshacl} 0.31.0 + \texttt{rdflib}
7.6.0 to the populated graph + composed shapes graph reproduces the
validation outcome: \texttt{sh:conforms\ =\ true}; zero
\texttt{sh:Violation}. The bounded-subset validator's parity for the six
ontology shapes is verified at the bundle's release time and preserved
at the bundle's manifest.

\subsubsection{11.4 Bundle artefact
identifiers}\label{bundle-artefact-identifiers}

The bundle's coordinated release uses the following stable artefact
identifiers, each SHA-256-pinned at the manifest:

\begin{itemize}
\tightlist
\item
  \texttt{ontology-v1-final} --- the OWL 2 DL ontology export with
  populated graph (§11.1);
\item
  \texttt{apparatus-shacl-pyshacl-v3} --- the SHACL apparatus shape
  files + build script (§11.2);
\item
  \texttt{appsec-core-embeddings-v1.1} --- the cycle's
  embedding-similarity tooling, used by the cycle reported in {[}4{]}
  and reproducible from the published embedding model and augmentation
  rule v1.0. The deposited embedding tensor is a \textbf{384-dimensional
  \texttt{float32} L2-normalized NPZ tensor} over a \textbf{212-entity
  augmented text corpus} (the four populated entity types'
  identifier-keyed text representations) under the encoder
  \texttt{sentence-transformers/all-MiniLM-L6-v2} {[}34{]} (corpus
  SHA-256 \texttt{5951fd82e4b7…}; embedding tensor SHA-256
  \texttt{17f6aac496b9…}; full 64-char SHA-256 digests recorded in the
  bundle manifest at the cycle-close release tag).
\end{itemize}

The bundle's audit-trail surface records the per-identifier introduction
provenance for each of the 25 net additions reported in §6 (Tables
2--3), with one verification entry per identifier and one
cross-reference per relation declaration.

\subsubsection{11.5 Per-instance attribution and release
notes}\label{per-instance-attribution-and-release-notes}

The bundle's release notes document per-instance attribution for each of
the four ACR worked decisions (ACR-001, ACR-002, ACR-004 promotions and
ACR-003 non-admission) and for the within-slice refinement additions
(§6.3). A reader auditing the v0→v1 transition against {[}1{]}'s
published surface can reconstruct the +25 net delta (§6.2) from the
per-decision evidence in §§8 and 9 plus the per-identifier introduction
record in the audit-trail surface (§11.4).

\begin{center}\rule{0.5\linewidth}{0.5pt}\end{center}

\subsection{AI Use Statement}\label{ai-use-statement}

Generative-language-model tools were used by the author during the
preparation of this paper for drafting and editing of prose, under the
author's review at every step. All scientific claims, methodological
decisions, evidence interpretation, and the per-identifier attribution
of the v0→v1 schema-preservation evidence (§6) are the author's
responsibility. The artefact this paper publishes --- AppSec Core v1,
its OWL 2 DL ontology export (§3), and the SHACL apparatus (§4) --- was
produced through the design science cycle reported in {[}4{]}, whose
substrate-generation pipeline invokes only a deterministic encoder
language model (Sentence-BERT-style sentence embeddings under
\texttt{sentence-transformers/all-MiniLM-L6-v2} {[}34{]}) at runtime; no
generative language model is invoked in the cycle's pipeline execution.
SHACL validation under \texttt{pyshacl} 0.31.0 + \texttt{rdflib} 7.6.0
(§4.2) is deterministic.

\begin{center}\rule{0.5\linewidth}{0.5pt}\end{center}

\subsection{12. References}\label{references}

\emph{References {[}4{]} and {[}5{]} are forthcoming programme papers;
at the time of this paper's deposit they are canonical via the
construction tags noted in the entry. OSF DOIs are assigned at
submission and supersede the construction-tag form.}

{[}1{]} P. Farinha. \emph{AppSec Core: A Normalized Ontology for
Security Requirements Across Heterogeneous Frameworks}. Independent
research preprint, 2026. DOI: 10.17605/OSF.IO/WG8PV.

{[}2{]} P. Farinha. \emph{Coverage-Preserving Compilation of Normative
and Empirical Security Knowledge}. Independent research preprint, 2026.
DOI: 10.17605/OSF.IO/A6ZFJ.

{[}3{]} P. Farinha. \emph{Research Programme Prospectus}. Independent
research preprint, 2026. DOI: 10.17605/OSF.IO/7T849.

{[}4{]} P. Farinha. \emph{Pressure-Testing AppSec Core: A Design Science
Cycle for Bounded-Ontology Evolution Under Heterogeneous
Application-Security Sources}. Independent research preprint (companion
to this paper), 2026. DOI: 10.17605/OSF.IO/3E8G5. Canonical via
construction tag v2.0.0-construction-p7-final-draft.

{[}5{]} P. Farinha. \emph{Coverage-Preserving Compilation v2: A
31-Source Pipeline with Joint Manual and Knowledge-Graph Production}.
Independent research preprint (succeeds {[}2{]} at 31-source scale),
2026 {[}in preparation; canonical via construction tag
v2.0.0-construction-p8-bundle-complete{]}.

{[}6{]} P. Farinha. \emph{Ontology-Grounded Retrieval for Auditable
LLM-Assisted Secure Code Generation}. Independent research preprint,
2026. DOI: 10.17605/OSF.IO/S3HET.

{[}7{]} P. Farinha. \emph{Empirical Evaluation: Deterministic Grounding
versus Retrieval-Augmented Generation versus Ungrounded Compilation}.
Pre-registered design, 2026 {[}empirical execution in preparation{]}.
DOI: 10.17605/OSF.IO/H5AJE.

{[}8{]} P. Farinha. \emph{MCP Dual-Mode Instrument Specification}.
Independent research preprint, 2026. DOI: 10.17605/OSF.IO/KH8Y7.

{[}9{]} National Institute of Standards and Technology. \emph{Secure
Software Development Framework (SSDF) Version 1.1}. NIST Special
Publication 800-218, 2022.

{[}10{]} OWASP Foundation. \emph{OWASP Application Security Verification
Standard (ASVS) Version 5.0}. 2025. URL:
https://owasp.org/www-project-application-security-verification-standard/.

{[}11{]} Open Source Security Foundation. \emph{Supply-chain Levels for
Software Artifacts (SLSA) Specification v1.0}. 2024. URL:
https://slsa.dev/spec/v1.0/.

{[}12{]} Center for Internet Security. \emph{CIS Critical Security
Controls Version 8.1.2}. 2024. URL:
https://www.cisecurity.org/controls/v8-1.

{[}13{]} MITRE Corporation. \emph{Common Attack Pattern Enumeration and
Classification (CAPEC), version 3.9}. URL: https://capec.mitre.org/,
accessed 2026.

{[}14{]} MITRE Corporation. \emph{Common Weakness Enumeration (CWE):
Software Development View, version 4.19.1}. URL: https://cwe.mitre.org/,
accessed 2026.

{[}15{]} OWASP Foundation. \emph{OWASP Guidance on Model Context
Protocol Security}. 2025 (referenced edition).

{[}16{]} OWASP Foundation. \emph{OWASP Software Assurance Maturity Model
(SAMM) Version 2.1}. 2023. URL: https://owaspsamm.org/.

{[}17{]} OWASP Foundation. \emph{OWASP DevSecOps Maturity Model
(DSOMM)}. 2024. URL: https://dsomm.owasp.org/.

{[}18{]} World Wide Web Consortium. \emph{Shapes Constraint Language
(SHACL)}. W3C Recommendation, 2017. URL: https://www.w3.org/TR/shacl/.

{[}19{]} World Wide Web Consortium. \emph{OWL 2 Web Ontology Language
Document Overview (Second Edition)}. W3C Recommendation, 2012. URL:
https://www.w3.org/TR/owl2-overview/.

{[}20{]} A. R. Hevner, S. T. March, J. Park, and S. Ram. ``Design
Science in Information Systems Research''. \emph{MIS Quarterly},
28(1):75--105, 2004.

{[}21{]} R. J. Wieringa. \emph{Design Science Methodology for
Information Systems and Software Engineering}. Springer, 2014.

{[}22{]} P. Ralph et al.~\emph{ACM SIGSOFT Empirical Standards}.
arXiv:2010.03525, 2021.

{[}23{]} N. Guarino. ``Formal Ontology and Information Systems''. In
\emph{Proceedings of FOIS 1998 --- 1st International Conference on
Formal Ontologies in Information Systems}, pp.~3--15. IOS Press, 1998.

{[}24{]} Z. Syed, A. Padia, T. Finin, L. Mathews, and A. Joshi. ``UCO: A
Unified Cybersecurity Ontology''. In \emph{AAAI Workshop on Artificial
Intelligence for Cyber Security}, 2016.

{[}25{]} M. D. Iannacone, S. Bohn, G. Nakamura, J. Gerth, K. Huffer, R.
A. Bridges, E. M. Ferragut, and J. R. Goodall. ``Developing an Ontology
for Cyber Security Knowledge Graphs''. In \emph{Proceedings of the 10th
Annual Cyber and Information Security Research Conference (CISRC)},
2015.

{[}26{]} C. Boyer, M. Dogdu, A. Choupani, R. Watson, A. Sanchez, and B.
Ametu. ``Toward a Unified Cybersecurity Knowledge Graph: Leveraging
Ontologies and Open Data Sources''. In \emph{IEEE International
Conference on Smart Data Services (SDSC)}, 2024.

{[}27{]} M. C. Suárez-Figueroa, A. Gómez-Pérez, and M. Fernández-López.
``The NeOn Methodology framework: A scenario-based methodology for
ontology development''. \emph{Applied Ontology}, 10(2):107--145, 2015.

{[}28{]} M. Poveda-Villalón, A. Fernández-Izquierdo, M. Fernández-López,
and R. García-Castro. ``LOT: An industrial oriented ontology engineering
framework''. \emph{Engineering Applications of Artificial Intelligence},
111:104755, 2022.

{[}29{]} K. Peffers, T. Tuunanen, M. A. Rothenberger, and S. Chatterjee.
``A Design Science Research Methodology for Information Systems
Research''. \emph{Journal of Management Information Systems},
24(3):45--77, 2007.

{[}30{]} B. Cuenca Grau, I. Horrocks, Y. Kazakov, and U. Sattler.
``Modular Reuse of Ontologies: Theory and Practice''. \emph{Journal of
Artificial Intelligence Research}, 31:273--318, 2008.

{[}31{]} B. Konev, C. Lutz, D. Walther, and F. Wolter. ``Formal
Properties of Modularisation''. In \emph{Modular Ontologies} (Springer
LNCS 5445), pp.~25--66, 2009.

{[}32{]} N. Keller, P. Quinn, and M. Smith. \emph{Mapping Relationships
Between Documentary Standards, Regulations, Frameworks, and Reference
Materials}. NIST Internal Report 8477, 2024.

{[}33{]} N. Keller, J. Barrett, P. Quinn, and M. Smith. \emph{National
Online Informative References (OLIR) Program: Submission Guidance for
OLIR Developers}. NIST Internal Report 8278A Revision 1, 2021.

{[}34{]} N. Reimers and I. Gurevych. ``Sentence-BERT: Sentence
Embeddings using Siamese BERT-Networks''. In \emph{Proceedings of
EMNLP-IJCNLP 2019}, pp.~3982--3992. arXiv:1908.10084.

{[}35{]} M. Poveda-Villalón, A. Gómez-Pérez, and M. C. Suárez-Figueroa.
``OOPS! (OntOlogy Pitfall Scanner!): An On-line Tool for Ontology
Evaluation''. \emph{International Journal on Semantic Web and
Information Systems}, 10(2):7--34, 2014. DOI: 10.4018/ijswis.2014040102.
URL: https://oops.linkeddata.es/.

{[}36{]} D. Garijo, O. Corcho, and M. Poveda-Villalón. ``FOOPS!: An
ontology pitfall scanner for the FAIR principles''. In \emph{Proceedings
of the ISWC 2021 Posters, Demos and Industry Tracks}, CEUR Workshop
Proceedings Vol. 2980, 2021. URL: https://foops.linkeddata.es/.

{[}37{]} National Institute of Standards and Technology. \emph{Security
and Privacy Controls for Information Systems and Organizations}. NIST
Special Publication 800-53 Revision 5, 2020. DOI:
10.6028/NIST.SP.800-53r5.

{[}38{]} PCI Security Standards Council. \emph{Payment Card Industry
Data Security Standard (PCI DSS) Version 4.0.1}, June 2024. URL:
https://www.pcisecuritystandards.org/document\_library/.

{[}39{]} PCI Security Standards Council. \emph{Payment Card Industry
Software Security Standard --- Secure Software Lifecycle (PCI SSLC)
Requirements Version 1.1}, 2022. URL:
https://www.pcisecuritystandards.org/document\_library/.

{[}40{]} European Parliament and Council. \emph{Regulation (EU)
2024/2847 on horizontal cybersecurity requirements for products with
digital elements (Cyber Resilience Act, CRA)}, October 2024. URL:
https://eur-lex.europa.eu/eli/reg/2024/2847.

{[}41{]} European Parliament and Council. \emph{Directive (EU) 2022/2555
on measures for a high common level of cybersecurity across the Union
(NIS2 Directive)}, December 2022. URL:
https://eur-lex.europa.eu/eli/dir/2022/2555.

{[}42{]} European Parliament and Council. \emph{Regulation (EU) 2016/679
(General Data Protection Regulation, GDPR)}, April 2016. URL:
https://eur-lex.europa.eu/eli/reg/2016/679.

{[}43{]} European Parliament and Council. \emph{Regulation (EU)
2022/2554 on digital operational resilience for the financial sector
(DORA)}, December 2022. URL:
https://eur-lex.europa.eu/eli/reg/2022/2554.

{[}44{]} U.S. Department of Health and Human Services. \emph{Health
Insurance Portability and Accountability Act (HIPAA) Security Rule}, 45
CFR Parts 160 and 164, Subparts A and C. URL:
https://www.hhs.gov/hipaa/for-professionals/security/.

{[}45{]} National Institute of Standards and Technology.
\emph{Adversarial Machine Learning: A Taxonomy and Terminology of
Attacks and Mitigations}. NIST AI 100-2 E2025, 2025. DOI:
10.6028/NIST.AI.100-2e2025.

{[}46{]} National Institute of Standards and Technology.
\emph{Artificial Intelligence Risk Management Framework (AI RMF 1.0)}.
NIST AI 100-1, January 2023. DOI: 10.6028/NIST.AI.100-1.

{[}47{]} OWASP Foundation. \emph{OWASP Proactive Controls}, 2018
edition. URL: https://owasp.org/www-project-proactive-controls/.

{[}48{]} OWASP Foundation. \emph{OWASP Top 10 for Large Language Model
Applications (LLM01--LLM10)}, 2025. URL:
https://genai.owasp.org/llm-top-10/.

{[}49{]} MITRE Corporation. \emph{Adversarial Threat Landscape for
Artificial-Intelligence Systems (ATLAS) v5.6.0}. URL:
https://atlas.mitre.org/.

{[}50{]} SAFECode. \emph{Practices for Secure Development of Cloud
Applications; Fundamental Practices for Secure Software Development;
Software Integrity Controls} (SAFECode Agile 2012, FPSSD 2018, SIC
2010). URL: https://safecode.org/publications/.

\end{document}
